% AUTO-GENERATED by Github-Branch/tools/build_unified_source.ps1.
% Do not edit this file directly: update the transformation manifest instead.
\chapter{Affine Stochastic Models with Stochastic Volatility and Jumps}
\chapterstatus{Edited reuse from Team 6. Synthetic illustrations and calibration claims are excluded; formula parity and branch conventions remain under review.}
\sourceassetnote
\section{From Conditional Variance to a Continuous-Time Pricing Baseline}

Chapter~7 established a discrete-time empirical fact: the conditional variance
of gold returns is dynamic and persistent even when the conditional mean is
weak. This chapter changes purpose as well as time scale. Its task is not to
replace GARCH as an econometric diagnostic, but to define risk-neutral
continuous-time processes that can price an option surface and later generate
internally consistent scenario paths.

Black--Scholes--Merton is the baseline. For a European claim
$C=C(S,t)$ with constant volatility $\sigma$, the replication argument gives
\[
 C_t+\frac{1}{2}\sigma^2S^2C_{SS}+rSC_S-rC=0.
\]
After changing from calendar time to time-to-maturity, taking logarithmic
moneyness as the spatial variable and rescaling the dependent variable, this
parabolic equation reduces to the standard heat equation. The connection is a
mathematical transformation that makes diffusion tools available to finance;
it is not a claim that Black--Scholes or Heston were generically ``taken from
physics.'' The same baseline also makes the later departures precise: constant
variance is replaced by a state process, and continuous paths are augmented by
jumps.

\section{Limits of the Classical Black–Scholes Framework}
\subsection{Continuous Paths and the Absence of Jumps}

The Black-Scholes-Merton framework fundamentally assumes that the underlying asset price follows a Geometric Brownian Motion (GBM). Under the risk-neutral measure $\mathbb{Q}$, the Stochastic Differential Equation (SDE) governing the asset is:
$$dS_t = r S_t dt + \sigma S_t dW_t^\mathbb{Q}$$
where $r$ is the constant risk-free rate, $\sigma$ is the constant volatility, and $W_t^\mathbb{Q}$ is a standard Brownian motion.A defining mathematical property of standard Brownian motion is that its sample paths are almost surely continuous. Because the only source of randomness in the GBM equation is the infinitesimal increment $dW_t^\mathbb{Q}$, the asset price process $S_t$ strictly inherits this continuity.We can formalize this by defining the left limit of the price process at time $t$ as $S_{t^-} = \lim_{s \to t^-} S_s$. The continuous-path assumption rigorously implies that the price immediately before the instant $t$ and the price at the exact instant $t$ are identical. Consequently, the instantaneous price jump, denoted as $\Delta S_t$, is strictly zero:
$$\Delta S_t = S_t - S_{t^-} = 0 \quad \text{with probability 1}$$
\subsubsection{The Empirical Contradiction and Delta Hedging Failure}

While mathematically elegant, this assumption strongly contradicts the reality of financial markets. Real-world assets frequently experience abrupt price discontinuities—or "jumps"—caused by overnight market closures, macroeconomic announcements, or sudden corporate defaults. In these practical scenarios, the jump size is not zero ($\Delta S_t \neq 0$).This limitation is critical because the entire Black-Scholes pricing logic relies on Delta hedging to continuously eliminate risk. To construct a perfectly risk-free portfolio, a trader must be able to continuously rebalance the underlying asset as its price smoothly diffuses. If an asset price gaps instantaneously from $100$ to $80$, the trader physically and mathematically cannot trade at the intermediate prices (e.g., $95$, $90$).When a jump occurs, the first-order Taylor expansion (Itô's Lemma) used to hedge the option breaks down. The continuous dynamic replication fails, making the market "incomplete" and leaving the hedged portfolio exposed to severe directional risk. Therefore, to safely price and hedge in real-world conditions, the framework must be expanded beyond pure diffusions to incorporate discontinuous stochastic processes.
\subsection{Constant Volatility and the Volatility Smile}
A core structural assumption of the Black-Scholes-Merton model is that the instantaneous volatility of the underlying asset returns, denoted by $\sigma$, is a deterministic constant over the life of the derivative. Under this assumption, the theoretical price of a European Call option $C_{BS}$ is given by the closed-form solution:
$$C_{BS}(S_t, K, \tau, \sigma) = S_t N(d_1) - K e^{-r\tau} N(d_2)$$
where $\tau = T - t$ is the time to maturity, and $d_1, d_2$ are defined as:
$$d_1 = \frac{\ln(S_t / K) + (r + \frac{1}{2}\sigma^2)\tau}{\sigma \sqrt{\tau}}, \quad d_2 = d_1 - \sigma \sqrt{\tau}$$
Because the partial derivative of the option price with respect to volatility—known as Vega ($\mathcal{V}$)—is strictly positive everywhere:$$\mathcal{V} = \frac{\partial C_{BS}}{\partial \sigma} = S_t \sqrt{\tau} \phi(d_1) > 0$$the Black-Scholes pricing function is strictly monotonically increasing with respect to $\sigma$. This mathematical property guarantees that for any observed, arbitrage-free market price of an option $C_{mkt}$, there exists a unique positive root $\sigma_{imp}$ that equates the theoretical model to the market reality. This root is defined as the Implied Volatility:
$$C_{mkt}(K, \tau) = C_{BS}(S_t, K, \tau, \sigma_{imp})$$
If the foundational assumption $d\sigma = 0$ were empirically correct, solving this inverse problem across a spectrum of different strike prices $K$ should yield a constant value for $\sigma_{imp}$. Geometrically, plotting $\sigma_{imp}$ as a function of $K$ would result in a perfectly flat horizontal line.However, empirical market data systematically violates this condition. By extracting implied volatilities across various strikes for a given maturity, one observes that:
$$\frac{\partial \sigma_{imp}}{\partial K} \neq 0 \quad \text{and} \quad \frac{\partial^2 \sigma_{imp}}{\partial K^2} > 0$$
This non-linear relationship forms a convex curve commonly referred to as the volatility smile or, more typically in equity markets, a pronounced downward-sloping volatility skew. Deep out-of-the-money (OTM) put options command significantly higher implied volatilities than at-the-money (ATM) options. This structural deformation occurs because market participants inherently price in the "fat tails" (excess kurtosis) and negative skewness of the real-world return distributions. To reconcile the mathematical model with observed market dynamics, the classical framework must be abandoned in favor of models where volatility itself follows a stochastic diffusion process.

% Asset/code block omitted pending provenance acceptance.
\subsection{Motivation for Jump and Stochastic Volatility Models}
Given the empirical failures of continuous sample paths and constant volatility assumptions, the classical Black-Scholes-Merton framework must be expanded. To accurately capture both the short-term and long-term topological features of the implied volatility surface, modern quantitative finance introduces two distinct but complementary mathematical extensions: Stochastic Volatility (SV) and Jump-Diffusion (JD) processes.
\subsubsection{The Case for Stochastic Volatility (Long-Maturity Skew)}
To capture the empirical observation that volatility fluctuates and clusters over time, SV models relax the assumption that $d\sigma = 0$. Instead, the variance $v_t = \sigma_t^2$ is promoted to a strictly positive state variable governed by its own stochastic differential equation, typically a mean-reverting square-root diffusion such as the Cox-Ingersoll-Ross (CIR) process. The coupled system under the risk-neutral measure takes the form:
$$dS_t = r S_t dt + \sqrt{v_t} S_t dW_t^S$$
$$dv_t = \kappa(\theta - v_t)dt + \xi \sqrt{v_t} dW_t^v$$

where the Brownian motions $dW_t^S$ and $dW_t^v$ are correlated ($d\langle W^S, W^v \rangle_t = \rho dt$). The correlation parameter $\rho < 0$ mathematically generates the "leverage effect," naturally producing a downward-sloping volatility skew. However, because SV diffusions are continuous, the variance requires time to accumulate. As time to maturity approaches zero ($\tau \to 0$), the diffusion variance $\int_t^T v_s ds$ converges to zero, causing the theoretical short-term volatility skew to flatten out, which contradicts the steep short-term smiles observed in equity markets.
\subsubsection{The Case for Jump Processes (Short-Maturity Skew)}
To address short-maturity pricing effects and the possibility of sudden market
repricing, jump--diffusion models augment the continuous Brownian component with
a jump component driven by a Poisson process $N_t$. The asset dynamics are
modified as
\[
    \frac{dS_t}{S_{t^-}}
    =
    (r-\lambda k)\,dt+\sigma\,dW_t+(e^J-1)\,dN_t,
    \qquad
    k:=\mathbb{E}^{\mathbb{Q}}[e^J-1],
\]
where $\lambda$ is the jump intensity and $J$ is the logarithmic jump size. The
compensator $\lambda k$ ensures that the discounted price is a martingale under
the pricing measure. Jumps can generate heavy tails and skewness even over short
time horizons, because a discontinuous move may occur before the diffusion
component has accumulated substantial variance. Over longer horizons, however,
the standardized contribution of finite-variance independent jumps becomes less
pronounced: the total variance grows linearly with maturity, while skewness and
excess kurtosis scale down. As a result, pure jump--diffusion models typically
struggle to reproduce persistent long-maturity volatility skews.


\textbf{Synthesis}
Ultimately, neither framework is independently sufficient. Jump processes are mathematically essential to capture abrupt discontinuities and price short-maturity out-of-the-money options, while stochastic volatility is required to model the persistent, mean-reverting nature of market variance and price long-maturity options. This duality motivates the transition toward unified affine jump-diffusion frameworks, such as the Bates model, which integrate both mechanisms to consistently describe the entire volatility surface.

% ---- next approved source slice ----

\section{Poisson Processes and Jump Times}
\subsection{Counting Processes}
To overcome the strict continuity assumptions of the classical diffusion framework, we must introduce mathematical structures capable of modeling discrete, isolated events occurring randomly in continuous time. The foundational building block for such models is the counting process.
Let $(\Omega, \mathcal{F}, \mathbb{P})$ be a probability space equipped with a filtration $\{\mathcal{F}_t\}_{t \ge 0}$ satisfying the usual conditions. A stochastic process $\{N_t\}_{t \ge 0}$ is defined as a counting process if it represents the total number of specific events that have occurred up to and including time $t$.
Mathematically, $\{N_t\}_{t \ge 0}$ must satisfy the following fundamental properties:
\begin{enumerate}
    \item $N_0 = 0$ almost surely
    \item For any time $t \ge 0$, $N_t$ takes values in the set of non-negative integers $\mathbb{N}_0 = \{0, 1, 2, \dots\}$
    \item The paths of the process are non-decreasing: for any two time points $s < t$, we have $N_s \le N_t$ almost surely.
    \item For any $s < t$, the increment $(N_t - N_s)$ represents the strictly non-negative integer number of events occurring in the time interval $(s, t]$.
\end{enumerate}
A crucial topological feature of a counting process is its sample path trajectory. Unlike standard Brownian motion $W_t$, which possesses continuous sample paths, the counting process $N_t$ is a piecewise constant step function. Its paths belong to the space of càdlàg functions (right-continuous with left limits).
If we denote the random arrival time of the $n$-th event as $\tau_n$, where $0 < \tau_1 < \tau_2 < \dots < \tau_n < \dots$, the instantaneous behavior of the process at an arrival time exhibits a discrete jump:
$$\Delta N_{\tau_n} = N_{\tau_n} - N_{\tau_n^-} = 1$$
where $N_{\tau_n^-} = \lim_{s \to \tau_n^-} N_s$. 

At all other times $t \neq \tau_n$, the differential $dN_t$ is strictly zero.Consequently, the counting process provides the exact mathematical language required to formalize the arrival of abrupt price discontinuities in financial markets, completely independent of the continuous diffusion component.

% Asset/code block omitted pending provenance acceptance.
\subsection{Homogeneous Poisson Processes}
While a general counting process provides the structural framework for discrete jumps, we must specify the probabilistic laws governing their arrival. The most fundamental model for random, independent jump arrivals in continuous time is the Homogeneous Poisson Process.
Given the filtered probability space $(\Omega, \mathcal{F}, \mathbb{P})$, a counting process $\{N_t\}_{t \ge 0}$ is defined as a Homogeneous Poisson Process with constant intensity (or rate) parameter $\lambda > 0$ if it satisfies the following axiomatic properties:
\begin{enumerate}
    \item \textbf{Independent Increments:} For any sequence of time points $0 \le t_0 < t_1 < \dots < t_n$, the random variables representing the number of events in each interval, $(N_{t_1} - N_{t_0}), (N_{t_2} - N_{t_1}), \dots, (N_{t_n} - N_{t_{n-1}})$, are mutually independent. This implies that the probability of a future jump is strictly independent of the past filtration $\mathcal{F}_s$ for $s < t$.
    \item \textbf{Stationary Increments:} The distribution of the number of events in any time interval depends exclusively on the length of the interval, and not on its absolute position in time.
    \item \textbf{Infinitesimal Transition Probabilities:} In an infinitesimally small time interval $dt$, the probability of observing exactly one jump is proportional to $\lambda$, while the probability of observing multiple simultaneous jumps is negligible. Formally:
    $$\mathbb{P}[N_{t+dt} - N_t = 1] = \lambda dt + o(dt)$$
    $$\mathbb{P}[N_{t+dt} - N_t \ge 2] = o(dt)$$
    $$\mathbb{P}[N_{t+dt} - N_t = 0] = 1 - \lambda dt + o(dt)$$
\end{enumerate}

A direct mathematical consequence of these axioms is that for any finite time interval of length $\tau = t - s > 0$, the increment $(N_t - N_s)$ follows a Poisson distribution with parameter $\lambda \tau$. The probability mass function (PMF) of observing exactly $k$ jumps in this interval is given by:
$$\mathbb{P}[N_t - N_s = k] = \frac{e^{-\lambda \tau} (\lambda \tau)^k}{k!}, \quad \text{for } k = 0, 1, 2, \dots$$
This specific distribution perfectly defines the moments of the process. For a Homogeneous Poisson Process, both the expected value and the variance grow linearly with time, scaled by the intensity $\lambda$:
$$\mathbb{E}[N_t] = Var(N_t) = \lambda t$$
In the context of quantitative finance, $\lambda$ represents the mean arrival rate of market shocks per unit of time (e.g., jumps per year). By assigning this strict probabilistic structure, we can mathematically price the probability of observing $k$ extreme market events before an option's maturity date.
\subsection{Waiting Times and Exponential Interarrival Structure}
The Homogeneous Poisson Process can be equivalently defined through the continuous time intervals elapsed between consecutive jumps. Rather than counting the number of events in a fixed time window, we analyze the random time required for the next event to occur.Let $\tau_1, \tau_2, \dots, \tau_n$ be the strictly increasing random arrival times of the jumps, where $\tau_1$ is the time of the first event. We define the sequence of interarrival times (or waiting times) $\{T_n\}_{n \ge 1}$ as:
$$T_1 = \tau_1$$
$$T_n = \tau_n - \tau_{n-1}, \quad \text{for } n \ge 2$$
For a Homogeneous Poisson Process with intensity $\lambda > 0$, the interarrival times $T_1, T_2, \dots$ form a sequence of independent and identically distributed (i.i.d.) random variables following an Exponential distribution.The probability density function (PDF) of each interarrival time $T_n$ is:
$$f_T(t) = \lambda e^{-\lambda t}, \quad \text{for } t \ge 0$$
Integrating this density, we can determine the probability that the waiting time exceeds a certain threshold $t$. In reliability engineering and mathematical finance, this is known as the survival probability, formalized as $\mathbb{P}[X > x] = 1 - F_X(x)$. For our interarrival times, the survival probability assumes the explicit form:
$$\mathbb{P}[T_n > t] = \int_t^\infty \lambda e^{-\lambda s} ds = e^{-\lambda t}$$
The expected waiting time between two consecutive jumps is the reciprocal of the intensity parameter:
$$\mathbb{E}[T_n] = \frac{1}{\lambda}$$
This provides a highly intuitive financial metric: if a specific market shock is modeled with an intensity of $\lambda = 0.5$ events per year, the expected waiting time between such shocks is $2$ years.
\paragraph{The Memoryless Property}~\\
A profound structural consequence of the exponential interarrival times is the memoryless property. Mathematically, for any $s, t \ge 0$:
$$\mathbb{P}[T_n > s + t \mid T_n > s] = \mathbb{P}[T_n > t]$$
In the context of financial markets, this property implies that the homogeneous jump process possesses no memory of the past. If a market has survived for a time $s$ without experiencing a crash, the conditional probability of surviving for an additional time $t$ depends solely on $t$ and $\lambda$. The risk of a jump arrival does not deterministically "age" or accumulate over time; the probability of an instantaneous shock remains perpetually constant.
\subsection{Financial Interpretation of Random Jump Arrivals}
The rigorous mathematical framework of the Poisson process provides a natural language for describing the mechanics of information flow in financial markets. In classical diffusion models, asset prices evolve continuously, which implicitly assumes that new information arrives in a smooth, continuous stream. While this adequately models normal daily trading noise and minor order-book imbalances, it fails to capture the true nature of critical macroeconomic and firm-specific news.

In reality, highly impactful information arrives discretely. The random jump times $\tau_1, \tau_2, \dots$ generated by the Poisson process map directly to the exact moments these sudden information shocks hit the market. We can broadly categorize these random arrivals into two distinct financial classes:
\begin{enumerate}
    \item \textbf{Idiosyncratic (Firm-Specific) Shocks:} These are localized events affecting a single asset, such as an unexpected earnings announcement, a sudden CEO resignation, a merger and acquisition (M\&A) bid, or a corporate default.
    \item \textbf{Systemic (Macroeconomic) Shocks:} These represent broader market disruptions affecting entire sectors or indices, such as surprise interest rate changes by central banks, geopolitical conflicts, or sudden liquidity crises.
\end{enumerate}
The intensity parameter $\lambda$ serves as a quantitative measure of the market's "hazard rate." A low $\lambda$ characterizes a stable economic environment with rare discontinuities, whereas a high $\lambda$ reflects a volatile, distressed market regime where sudden shocks are expected to arrive with higher frequency.
\paragraph{The Bridge to Compound Processes}~\\
While the simple counting process $N_t$ successfully models when a market shock occurs (yielding a jump of magnitude $\Delta N_t = 1$), it does not quantify the economic severity of the shock. An interest rate hike might cause a $2\%$ drop in an equity index, whereas a corporate bankruptcy could erase $80\%$ of a stock's value instantaneously.

To build a complete asset pricing model, the fixed unit jump of the counting process must be replaced with a random variable representing the jump size. This mathematical evolution—combining the random arrival times of the Poisson process with random price impacts—leads directly to the formulation of the Compound Poisson Process and Jump-Diffusion models.
\section{Compound Poisson Processes and Jump--Diffusions}
\subsection{Random Jump Sizes}

The Poisson process introduced previously provides a natural description of the random arrival times of events. In order to construct processes with jumps, it is necessary to associate to each jump time a random amplitude.

Let $(\Omega, \mathcal{F}, \mathbb{P})$ be a probability space and let $\{N_t\}_{t \geq 0}$ be a Poisson process with intensity $\lambda > 0$. We consider a sequence of random variables $\{Y_i\}_{i \geq 1}$ defined on the same probability space, such that:

\begin{itemize}[label = --]
    \item the random variables $Y_i$ are independent and identically distributed;
    \item the sequence $\{Y_i\}$ is independent of the process $\{N_t\}_{t \geq 0}$.
\end{itemize}

The random variables $Y_i$ represent the jump sizes and take values in $\mathbb{R}$. This construction allows one to associate to each jump time $\tau_i$ of the Poisson process a mark $Y_i$, thus obtaining a marked point process.

From a modelling perspective, the independence assumption reflects the fact that the arrival times of jumps and their magnitudes are driven by distinct sources of randomness.

\medskip

This framework naturally leads to the construction of jump processes by summing the contributions of individual jumps.
\subsection{Compound Poisson Sums}

\begin{definition}
The stochastic process $\{X_t\}_{t \geq 0}$ defined by
\[
X_t = \sum_{i=1}^{N_t} Y_i
\]
is called a \emph{compound Poisson process}.
\end{definition}

By convention, the sum is taken to be zero when $N_t = 0$.

Conditionally on the event $\{N_t = n\}$, the random variable $X_t$ is equal to the sum of $n$ independent random variables with the same distribution as $Y_1$. Therefore, its distribution is given by the $n$-fold convolution of the law of $Y_1$.

Using the law of total expectation, one obtains:
\[
\mathbb{E}[X_t] = \mathbb{E}[N_t] \mathbb{E}[Y_1] = \lambda t \, \mathbb{E}[Y_1],
\]
and
\[
\mathrm{Var}(X_t) = \lambda t \, \mathbb{E}[Y_1^2].
\]

The process $\{X_t\}_{t \geq 0}$ has independent and stationary increments and is therefore a L\'evy process. Its sample paths are piecewise constant, right-continuous with left limits (c\`adl\`ag), with jumps occurring at the jump times of the Poisson process.

\medskip

An equivalent representation is obtained by introducing the discrete-time random walk
\[
R(n) = \sum_{i=1}^{n} Y_i, \quad n \in \mathbb{N},
\]
so that
\[
X_t = R(N_t).
\]

This representation emphasizes that a compound Poisson process can be viewed as a random walk observed at random times.

\medskip

A fundamental property of compound Poisson processes is that they belong to the class of L\'evy processes with \emph{finite activity}, meaning that the number of jumps on any finite time interval is almost surely finite.

More precisely, for any $t \geq 0$, the number of jumps of $X_t$ on $[0,t]$ is equal to $N_t$, which is finite with probability one.

\medskip

Finally, the characteristic function of $X_t$ is given by
\[
\mathbb{E}\left[e^{iu X_t}\right] = \exp\left( \lambda t \left( \mathbb{E}[e^{iuY_1}] - 1 \right) \right),
\]
which follows from conditioning on $N_t$ and using the independence of the jump sizes.
This expression highlights the infinite divisibility of the distribution of $X_t$ and plays a central role in the analysis of jump processes.
\subsection{From Diffusions to Jump--Diffusions}


The compound Poisson process captures the abrupt component of price
dynamics, but its sample paths are piecewise constant between jumps
and therefore cannot accommodate the continuous fluctuations that
Brownian motion describes so effectively. A realistic model of the
underlying asset price must superpose the two mechanisms: a diffusive
component for the ordinary trading noise and a compound Poisson
component for the discontinuous revaluations of the underlying
asset. This programme, originally proposed by
Merton, yields the class of \emph{jump--diffusion} models.

Let $(\Omega, \mathcal{F}, \{\mathcal{F}_t\}_{t\geq 0}, \mathbb{P})$
be a filtered probability space supporting
a standard Brownian motion $\{W_t\}_{t\geq 0}$, a Poisson process
$\{N_t\}_{t\geq 0}$ with intensity $\lambda > 0$ and a sequence
$\{Y_i\}_{i\geq 1}$ of i.i.d.\ real random variables with
$\mathbb{E}[e^{Y_1}] < \infty$. We assume that $W$, $N$ and
$\{Y_i\}_{i\geq 1}$ are mutually independent and that
$\{\mathcal{F}_t\}_{t\geq 0}$ is the filtration jointly generated by
them, so that $W$ and $N$ are Brownian and Poisson relative to this
filtration.
The asset price $\{S_t\}_{t\geq 0}$ is then postulated to evolve
according to the stochastic differential equation
\begin{equation}
    dS_t = \mu\, S_{t^-}\, dt + \sigma\, S_{t^-}\, dW_t
         + S_{t^-}\!\left(e^{Y} - 1\right) dN_t,
    \qquad S_0 > 0,
\end{equation}
with constants $\mu \in \mathbb{R}$, $\sigma > 0$ and $Y$ denoting the
generic element of the sequence $\{Y_i\}_{i\geq 1}$ that provides the
jump sizes at the successive arrival times of $N$. The left limit
$S_{t^-}$ appears in the Brownian and Poisson integrals as required
by the predictability of the integrand: between jumps
$S_{t^-} = S_t$, so the first two terms reproduce the usual geometric
Brownian dynamics, whereas at a jump time $\tau_i$ of $N$ the
increment of the third term contributes $S_{\tau_i^-}(e^{Y_i} - 1)$,
so that
\[
    S_{\tau_i} = S_{\tau_i^-}\, e^{Y_i}.
\]
The sample paths of $\{S_t\}_{t\geq 0}$ are therefore almost surely
c\`adl\`ag, strictly positive, and coincide with a geometric Brownian
motion on each stochastic interval between two consecutive arrival
times of $N$; at each arrival time they are multiplied by the
independent random factor $e^{Y_i}$.

An explicit solution of~\sourceref{eq:jumpdiff_SDE} can be obtained by
factoring $S_t$ into its continuous and pure-jump components and
applying the It\^o--Doeblin formula for jump processes. Define
\[
    X_t = S_0 \exp\!\left\{\sigma W_t
        + \left(\mu - \tfrac{1}{2}\sigma^2\right) t\right\},
    \qquad
    J_t = \prod_{i=1}^{N_t} e^{Y_i},
\]
with the empty product equal to $1$. The process $\{X_t\}_{t\geq 0}$
is continuous and satisfies, by the It\^o formula for continuous
semimartingales,
\[
    dX_t = \mu\, X_t\, dt + \sigma\, X_t\, dW_t,
\]
while $\{J_t\}_{t\geq 0}$ is a pure-jump process whose jump at the
arrival time $\tau_i$ of $N$ equals $J_{\tau_i^-}(e^{Y_i} - 1)$, so
that in differential form
\[
    dJ_t = J_{t^-}\!\left(e^{Y} - 1\right) dN_t.
\]
Because $X$ is continuous and $J$ is a pure-jump process, their
quadratic covariation vanishes. Applying the product rule for jump
processes to $S_t = X_t J_t$ yields
\[
    dS_t = X_{t^-}\, dJ_t + J_t\, dX_t
         = \mu\, S_{t^-}\, dt + \sigma\, S_{t^-}\, dW_t
         + S_{t^-}\!\left(e^{Y} - 1\right) dN_t,
\]
which is precisely~\sourceref{eq:jumpdiff_SDE}. We have therefore
established the closed-form representation
\begin{equation}
    S_t = S_0 \exp\!\left\{\sigma W_t
        + \left(\mu - \tfrac{1}{2}\sigma^2\right) t\right\}
        \prod_{i=1}^{N_t} e^{Y_i}
        = S_0 \exp\!\left\{\sigma W_t
            + \left(\mu - \tfrac{1}{2}\sigma^2\right) t
            + \sum_{i=1}^{N_t} Y_i\right\}.
\end{equation}
The logarithmic price $\ln(S_t/S_0)$ is thus the sum of a Brownian
component with drift $\mu - \tfrac{1}{2}\sigma^2$ and an independent
compound Poisson component $\sum_{i=1}^{N_t} Y_i$ with intensity
$\lambda$ and jump-size distribution that of $Y_1$. In differential
form,
\begin{equation}
    d\ln S_t = \left(\mu - \tfrac{1}{2}\sigma^2\right) dt
             + \sigma\, dW_t + Y\, dN_t,
\end{equation}
where $Y\,dN_t$ denotes the increment of the compound Poisson process
$\sum_{i=1}^{N_t} Y_i$ at a jump time of $N$.

% Asset/code block omitted pending provenance acceptance.


Representation~\sourceref{eq:jumpdiff_sol} makes transparent the two
structural features that distinguish jump--diffusions from the
Black--Scholes dynamics. First, the
distribution of $\ln(S_t/S_0)$ is no longer Gaussian but a mixture of
Gaussians indexed by $N_t$: conditional on $\{N_t = k\}$, the log
return is normal with mean
$\bigl(\mu - \tfrac{1}{2}\sigma^2\bigr) t + \sum_{i=1}^{k} Y_i$
and variance $\sigma^2 t$, and the unconditional law aggregates these
components weighted by the Poisson probabilities
$e^{-\lambda t}(\lambda t)^k / k!$. This mixture structure generates
the excess kurtosis and asymmetric skewness that the pure-diffusion
model failed to reproduce. Secondly, with probability
$1 - e^{-\lambda t} > 0$ the trajectory $s \mapsto S_s$ exhibits at
least one discontinuity on the interval $[0,t]$, in sharp contrast
with the almost-sure continuity of the Black--Scholes price process
on the same interval. The dynamics~\sourceref{eq:jumpdiff_SDE} thus
removes the two most consequential structural limitations of the
diffusive framework simultaneously, at the cost of introducing a
second, discontinuous source of randomness whose implications for
hedging and for the uniqueness of the equivalent martingale measure
will be analysed in the next subsection.
\subsection{Drift Compensation and Risk-Neutral Dynamics}

In an arbitrage-free setting, pricing is performed under a risk-neutral probability measure $\mathbb{Q}$, equivalent to the physical measure $\mathbb{P}$, under which discounted asset prices are martingales.

More precisely, if $\{S_t\}_{t \geq 0}$ denotes the asset price process, the no-arbitrage condition requires that the discounted process satisfies
\[
e^{-rt} S_t \quad \text{is a martingale under } \mathbb{Q}.
\]

Equivalently, for all $t \geq 0$,
\[
\mathbb{E}^{\mathbb{Q}}[S_t \mid \mathcal{F}_0] = S_0 e^{rt}.
\]

In the presence of jumps, this martingale condition imposes a constraint on the drift of the process. To make this explicit, we rewrite the jump--diffusion dynamics in differential form:
\[
\frac{dS_t}{S_{t^-}} = \mu \, dt + \sigma \, dW_t + (e^{Y} - 1) dN_t.
\]

Taking expectations under $\mathbb{Q}$ and using the independence between the Brownian and Poisson components, together with the fact that $\mathbb{E}[dW_t] = 0$ and $\mathbb{E}[dN_t] = \lambda dt$, we obtain
\[
\mathbb{E}^{\mathbb{Q}}\left[\frac{dS_t}{S_{t^-}}\right] = \mu \, dt + \lambda \, \mathbb{E}[e^{Y} - 1] \, dt.
\]

The martingale condition for the discounted price process requires that the expected return equals the risk-free rate, that is
\[
\mathbb{E}^{\mathbb{Q}}\left[\frac{dS_t}{S_{t^-}}\right] = r \, dt.
\]

Comparing the two expressions yields the drift restriction
\[
\mu = r - \lambda \mathbb{E}[e^{Y} - 1].
\]

This adjustment term compensates for the average contribution of jumps, ensuring that the discounted price process remains a martingale under $\mathbb{Q}$.

It is often convenient to express the jump contribution in compensated form.
When jump sizes are random, the compensation must be applied to the marked jump
measure rather than only to the scalar Poisson process. Equivalently, setting
\[
    k:=\mathbb{E}^{\mathbb{Q}}[e^J-1],
\]
the risk-neutral dynamics can be written as
\[
    \frac{dS_t}{S_{t^-}}
    =
    r\,dt+\sigma\,dW_t^{\mathbb{Q}}
    +
    \Big((e^J-1)\,dN_t-\lambda k\,dt\Big).
\]
The bracketed term is the compensated jump contribution and has zero conditional
expectation under $\mathbb{Q}$. In the language of marked point processes, this
can be written more compactly as
\[
    \frac{dS_t}{S_{t^-}}
    =
    r\,dt+\sigma\,dW_t^{\mathbb{Q}}
    +
    \int_{\mathbb{R}}(e^y-1)\,\widetilde{N}(dt,dy),
\]
where $\widetilde{N}(dt,dy)$ denotes the compensated Poisson random measure with
mark distribution equal to the law of the jump size $J$.


\medskip

From a financial perspective, this result highlights a key difference with respect to the classical Black--Scholes framework. While in pure diffusion models the martingale condition uniquely determines the drift, in the presence of jumps the market becomes incomplete, and the equivalent martingale measure is, in general, not unique.

% ---- next approved source slice ----

\section{L\'evy Processes and Discontinuous Dynamics}

The compound Poisson processes introduced in the previous section already provide a first example of stochastic dynamics with discontinuous sample paths. Their jump activity is, however, finite on every compact time interval. In many applications, especially in quantitative finance, one needs a broader class of processes capable of combining a continuous martingale component, a drift term, and a possibly much richer jump structure. This leads naturally to the class of L\'evy processes.

L\'evy processes play a central role in modern stochastic modelling because they preserve, in continuous time, the two structural properties that make random walks tractable: independence and stationarity of increments. They therefore constitute the canonical continuous-time extension of sums of i.i.d.\ random variables and provide the natural probabilistic framework for discontinuous asset returns.
\subsection{Definition and Main Properties}

We work on a filtered probability space
\[
(\Omega,\mathcal{F},\{\mathcal{F}_t\}_{t\geq 0},\mathbb{P}),
\]
and consider a stochastic process $\{X_t\}_{t\geq 0}$ with values in~$\mathbb{R}$.

\begin{definition}
A stochastic process $\{X_t\}_{t\geq 0}$ is called a \emph{L\'evy process} if the following properties hold:
\begin{itemize}[label = --]
    \item $X_0 = 0$ almost surely;
    \item $\{X_t\}_{t\geq 0}$ has \emph{independent increments}, namely for every choice of times
    \[
    0 \leq t_0 < t_1 < \dots < t_n,
    \]
    the random variables
    \[
    X_{t_1}-X_{t_0},\; X_{t_2}-X_{t_1},\; \dots,\; X_{t_n}-X_{t_{n-1}}
    \]
    are independent;
    \item $\{X_t\}_{t\geq 0}$ has \emph{stationary increments}, namely for every $s,t \geq 0$ the law of
    \[
    X_{t+s}-X_s
    \]
    depends only on $t$;
    \item $\{X_t\}_{t\geq 0}$ is \emph{stochastically continuous}, that is, for every $\varepsilon > 0$ and every $t \geq 0$,
    \[
    \lim_{h\to 0}\mathbb{P}\bigl(|X_{t+h}-X_t|>\varepsilon\bigr)=0;
    \]
    \item the sample paths are almost surely c\`adl\`ag.
\end{itemize}
\end{definition}

\begin{remark}
The c\`adl\`ag regularity of the sample paths can in fact be derived from the remaining four properties: any process satisfying the first four conditions of the definition admits a modification whose trajectories are almost surely c\`adl\`ag (see Sato, 1999, Theorem~11.5). We include it in the definition in order to fix the version of the process with which we work throughout.
\end{remark}

The definition isolates the continuous-time analogue of a random walk. Indeed, if one observes a L\'evy process only on a deterministic grid $0,\Delta,2\Delta,\dots$, then the increments
\[
X_{\Delta}-X_0,\quad X_{2\Delta}-X_{\Delta},\quad \dots,\quad X_{n\Delta}-X_{(n-1)\Delta}
\]
are independent and identically distributed by stationarity and independence of increments. Thus the discretely sampled process
\begin{equation}
S_n^{(\Delta)} := X_{n\Delta}
\end{equation}
is a random walk.

Conversely, a L\'evy process may be viewed as a continuous-time interpolation principle for random walks. This interpretation is fundamental, because it explains why L\'evy processes inherit many structural properties that are familiar from discrete-time probabilistic models.

\medskip

A first immediate consequence of the definition is the following.

\begin{theorem}[Increment identity]
Let $\{X_t\}_{t\geq 0}$ be a L\'evy process. Then for every $0\leq s<t$,
\begin{equation}
X_t-X_s \overset{d}{=} X_{t-s}.
\end{equation}
\end{theorem}

\begin{proof}
By stationarity of increments, the law of $X_t-X_s$ depends only on the time difference $t-s$. Evaluating this increment over the interval $[0,t-s]$ yields
\[
X_t-X_s \overset{d}{=} X_{t-s}-X_0.
\]
Since $X_0=0$ almost surely, the claim follows.
\end{proof}

This identity shows that the one-dimensional marginal laws of a L\'evy process are completely determined by the family of increment distributions.

\medskip

Another important property concerns infinite divisibility. Since for any integer $n\geq 1$ one may decompose the interval $[0,t]$ into $n$ subintervals of equal length, one has
\[
X_t
=
\sum_{k=1}^{n}\bigl(X_{kt/n}-X_{(k-1)t/n}\bigr),
\]
where the summands are independent and identically distributed, each with the same law as $X_{t/n}$. Hence the law of $X_t$ is infinitely divisible.

\begin{theorem}[Infinite divisibility of marginal laws]
For every $t\geq 0$, the distribution of $X_t$ is infinitely divisible.
\end{theorem}

\begin{proof}
Fix $t>0$ and $n\in\mathbb{N}$. Write
\[
X_t
=
\sum_{k=1}^{n}\bigl(X_{kt/n}-X_{(k-1)t/n}\bigr).
\]
By independence of increments, the $n$ summands are independent. By stationarity of increments, they are identically distributed. Therefore the law of $X_t$ is the $n$-fold convolution of the law of $X_{t/n}$, and since $n$ is arbitrary, the distribution is infinitely divisible.
\end{proof}

This fact is not secondary. It is one of the key structural fingerprints of the class: every L\'evy process generates an infinitely divisible family of marginal laws. The converse also holds: infinitely divisible distributions are precisely the distributions that arise at fixed times from L\'evy processes, and this constitutes one of the deep consequences of the L\'evy--Khintchine representation theorem developed below.

\begin{remark}
Every L\'evy process is a strong Markov process. This follows from the independence of increments: for any stopping time~$\tau$, the future evolution $\{X_{\tau+t}-X_\tau\}_{t\geq 0}$ is independent of $\mathcal{F}_\tau$ and has the same law as $\{X_t\}_{t\geq 0}$. The Markov property is of central importance in financial applications, as it underlies the partial integro-differential equations used for option pricing under L\'evy dynamics.
\end{remark}
\subsection{Brownian Motion and Compound Poisson as Special Cases}

The definition above contains, as special cases, both the continuous diffusive model and the finite-activity jump model developed previously. This is one of the main reasons why the L\'evy framework is so useful: it unifies in a single class several stochastic mechanisms that, at first sight, may appear rather different.
\subsubsection{Brownian Motion}

Let $\{W_t\}_{t\geq 0}$ be a standard Brownian motion. Then:
\begin{itemize}[label = --]
    \item $W_0=0$ almost surely;
    \item Brownian motion has independent increments;
    \item for every $s,t\geq 0$,
    \[
    W_{t+s}-W_s \sim \mathcal{N}(0,t),
    \]
    so the increments are stationary;
    \item Brownian motion is continuous, hence in particular c\`adl\`ag and stochastically continuous.
\end{itemize}
Therefore Brownian motion is a L\'evy process. More generally, any process of the form
\begin{equation}
X_t = \mu t + \sigma W_t,
\end{equation}
with $\mu\in\mathbb{R}$ and $\sigma\geq 0$, is a L\'evy process. This class represents the purely continuous Gaussian component inside the general L\'evy framework, and its characteristic function is
\begin{equation}
\mathbb{E}\bigl[e^{iuX_t}\bigr] = \exp\!\left(iu\mu t - \tfrac{1}{2}\sigma^2 u^2 t\right), \qquad u\in\mathbb{R}.
\end{equation}
\subsubsection{Compound Poisson Processes}

Let $\{N_t\}_{t\geq 0}$ be a Poisson process with intensity $\lambda>0$, and let $\{Y_i\}_{i\geq 1}$ be i.i.d.\ real-valued random variables, independent of $N$. Define
\begin{equation}
X_t = \sum_{i=1}^{N_t} Y_i.
\end{equation}
As established above, $\{X_t\}_{t\geq 0}$ is a compound Poisson process with characteristic function
\begin{equation}
\mathbb{E}\bigl[e^{iuX_t}\bigr] = \exp\!\left(\lambda t\bigl(\mathbb{E}[e^{iuY_1}]-1\bigr)\right), \qquad u\in\mathbb{R}.
\end{equation}

We now verify that it is a L\'evy process.

\begin{theorem}[Compound Poisson is a L\'evy process]
Every compound Poisson process is a L\'evy process.
\end{theorem}

\begin{proof}
We have $X_0=0$ almost surely because $N_0=0$. Fix
\[
0\leq t_0 < t_1 < \dots < t_n.
\]
Then
\[
X_{t_k}-X_{t_{k-1}}
=
\sum_{i=N_{t_{k-1}}+1}^{N_{t_k}} Y_i.
\]
Since the Poisson process has independent increments and the marks $\{Y_i\}$ are independent of $N$ and i.i.d., the random sums over disjoint time intervals are independent. Hence the process has independent increments.

Moreover, the law of
\[
X_{t+s}-X_s
=
\sum_{i=N_s+1}^{N_{t+s}} Y_i
\]
depends only on the distribution of $N_{t+s}-N_s$, which is Poisson with parameter $\lambda t$, and on the common law of the jump sizes. Therefore the increments are stationary.

The sample paths are piecewise constant with finitely many jumps on compact intervals, hence c\`adl\`ag. Stochastic continuity follows from the fact that
\[
\mathbb{P}(N_{t+h}-N_t\geq 1)\to 0
\qquad \text{as } h\to 0.
\]
Thus $\{X_t\}_{t\geq 0}$ is a L\'evy process.
\end{proof}

This result clarifies the hierarchy introduced so far:
\begin{itemize}[label = --]
    \item Brownian motion is the basic continuous L\'evy process;
    \item compound Poisson processes are the basic finite-activity jump L\'evy processes;
    \item general L\'evy processes extend both by allowing continuous Gaussian motion and a jump structure that may be either finite or infinite in activity.
\end{itemize}
\subsection{C\`adl\`ag Paths and Independent Increments}

The analytical content of the definition of L\'evy process lies in the interaction between two ideas: the pathwise property of being c\`adl\`ag, and the distributional property of having stationary independent increments.
\subsubsection{C\`adl\`ag trajectories}

A function $f:[0,\infty)\to\mathbb{R}$ is c\`adl\`ag if it is right-continuous and admits a finite left limit at every strictly positive time. Thus, for every $t>0$,
\[
f(t^-):=\lim_{s\uparrow t}f(s)
\]
exists, and
\[
\lim_{s\downarrow t}f(s)=f(t).
\]

If $\{X_t\}_{t\geq 0}$ is a L\'evy process, then for almost every $\omega\in\Omega$ the map
\[
t\longmapsto X_t(\omega)
\]
is c\`adl\`ag. This is the natural regularity class for jump models: it allows discontinuities, but only of jump type. Oscillatory pathologies are excluded.

At any jump time $t$, the jump of the process is defined by
\begin{equation}
\Delta X_t := X_t - X_{t^-}.
\end{equation}
If $\Delta X_t\neq 0$, the process experiences a discontinuity at time $t$.

\medskip

This notation is particularly convenient because it isolates the discontinuous contribution of the process. In the compound Poisson case, for instance, the jumps are exactly the marked amplitudes $Y_i$ occurring at the arrival times of the underlying Poisson process, consistently with the notation established above.


The second structural ingredient is \emph{independence of increments}  This means that if one observes the process over disjoint time intervals, then the corresponding variations do not influence one another probabilistically.

To make this precise, if
\[
0\leq t_0 < t_1 < \dots < t_n,
\]
then the family
\[
X_{t_1}-X_{t_0},\; X_{t_2}-X_{t_1},\; \dots,\; X_{t_n}-X_{t_{n-1}}
\]
is independent.

This property is essential in continuous-time finance because it formalises the idea that the future innovation of the process, once a present time is fixed, is probabilistically decoupled from the past. When combined with stationarity, it implies that the distribution of future increments depends only on the length of the future horizon, not on the calendar date at which the horizon begins.

\medskip

In particular, if $0\leq s<t<u<v$, then
\[
X_t-X_s
\quad\text{and}\quad
X_v-X_u
\]
are independent random variables. This gives the process a strong renewal structure. It is precisely this structure that makes characteristic functions and transform methods especially effective in the study of L\'evy dynamics.
\subsubsection{A fundamental exponential representation}

Let
\begin{equation}
\varphi_t(u):=\mathbb{E}\bigl[e^{iuX_t}\bigr],
\qquad u\in\mathbb{R}.
\end{equation}
Because the increments are stationary and independent, the characteristic function of $X_t$ must have an exponential dependence on time.

\begin{theorem}[Existence of the characteristic exponent]
For every L\'evy process $\{X_t\}_{t\geq 0}$ there exists a function $\psi:\mathbb{R}\to\mathbb{C}$ such that
\begin{equation}
\mathbb{E}\bigl[e^{iuX_t}\bigr]
=
e^{\,t\psi(u)},
\qquad t\geq 0,\; u\in\mathbb{R}.
\end{equation}
The function $\psi$ is called the \emph{characteristic exponent} of the process.
\end{theorem}

\begin{proof}
Fix $u\in\mathbb{R}$. For any $s,t\geq 0$,
\[
X_{t+s} = X_t + (X_{t+s}-X_t).
\]
Since $X_t$ and $X_{t+s}-X_t$ are independent, and the second increment has the same law as $X_s$, we obtain
\[
\varphi_{t+s}(u)
=
\mathbb{E}\bigl[e^{iuX_t}\bigr]
\,\mathbb{E}\bigl[e^{iu(X_{t+s}-X_t)}\bigr]
=
\varphi_t(u)\,\varphi_s(u).
\]
Thus $t\mapsto \varphi_t(u)$ satisfies the multiplicative Cauchy equation on $\mathbb{R}_+$.

It remains to show that this functional equation, together with continuity, implies the exponential form. The stochastic continuity of $\{X_t\}_{t\geq 0}$ guarantees that $X_{t+h}\to X_t$ in probability as $h\to 0$. Convergence in probability implies convergence in distribution, which in turn is equivalent to pointwise convergence of characteristic functions (by L\'evy's continuity theorem). Hence the map $t\mapsto \varphi_t(u)$ is continuous.

A continuous solution $f:\mathbb{R}_+\to\mathbb{C}\setminus\{0\}$ of the multiplicative Cauchy equation $f(s+t)=f(s)f(t)$ with $f(0)=1$ is necessarily of the form $f(t)=e^{ct}$ for some $c\in\mathbb{C}$ (see Lukacs, 1970, Chapter~7). Setting $c=\psi(u)$ yields
\[
\varphi_t(u)=e^{t\psi(u)},
\]
which proves the claim.
\end{proof}

This representation is fundamental. It shows that once the exponent $\psi$ is known, the whole family of marginal characteristic functions is known. Therefore the probabilistic structure of a L\'evy process is encoded in a single deterministic function.

\medskip

\begin{remark}
For the two special cases treated above, the characteristic exponent takes the following explicit forms:
\begin{itemize}[label = --]
    \item Brownian motion with drift \sourceref{eq:BM_levy}:
    \[
    \psi(u) = i\mu u - \tfrac{1}{2}\sigma^2 u^2;
    \]
    \item compound Poisson process \sourceref{eq:CPP_def} with jump-size distribution $F$:
    \[
    \psi(u) = \lambda\bigl(\mathbb{E}[e^{iuY_1}]-1\bigr).
    \]
\end{itemize}
These expressions are consistent with the characteristic functions \sourceref{eq:BM_charfun} and \sourceref{eq:CPP_charfun} derived above.
\end{remark}
\subsection{The L\'evy--Khintchine Representation}


The preceding section shows that every L\'evy process admits a characteristic exponent $\psi$. The next question is structural: which functions $\psi$ can arise from a L\'evy process? The answer is given by the L\'evy--Khintchine formula, which provides a complete parametric characterisation of all admissible exponents.

\begin{theorem}[L\'evy--Khintchine formula]
A function $\psi:\mathbb{R}\to\mathbb{C}$ is the characteristic exponent of a L\'evy process if and only if it can be written in the form
\begin{equation}
\psi(u)
=
i\gamma u
-\frac{1}{2}Au^2
+
\int_{\mathbb{R}\setminus\{0\}}
\left(
e^{iux}-1-iux\,\mathbf{1}_{\{|x|\leq 1\}}
\right)\nu(dx),
\end{equation}
where:
\begin{itemize}[label = --]
    \item $\gamma\in\mathbb{R}$ is a drift parameter;
    \item $A\geq 0$ is the Gaussian variance coefficient;
    \item $\nu$ is a measure on $\mathbb{R}\setminus\{0\}$ satisfying the integrability condition
    \begin{equation}
    \int_{\mathbb{R}\setminus\{0\}}\min\{1,x^2\}\,\nu(dx)<\infty.
    \end{equation}
\end{itemize}
The measure $\nu$ is called the \emph{L\'evy measure}, and the triplet $(\gamma,A,\nu)$ is called the \emph{L\'evy triplet} (or \emph{characteristic triplet}) of the process.
\end{theorem}

The proof of the theorem is technically demanding and goes beyond the scope of the present work; for a complete treatment the reader is referred to Sato (1999, Theorem~8.1) and Applebaum (2009, Theorem~1.3.3). In the remainder of this section we focus on the interpretation of the formula and on the verification of its consistency with the special cases already studied.

The formula~\sourceref{eq:LevyKhintchine} should be read as a structural decomposition of all possible continuous-time processes with stationary independent increments.

Each term in \sourceref{eq:LevyKhintchine} corresponds to a distinct dynamical mechanism.

\begin{itemize}[label = --]
    \item The term
    \[
    i\gamma u
    \]
    represents the deterministic linear drift.
    \item The term
    \[
    -\frac{1}{2}Au^2
    \]
    is the Gaussian component. It corresponds to a Brownian motion with variance~$A$.
    \item The integral term
    \[
    \int_{\mathbb{R}\setminus\{0\}}
    \left(
    e^{iux}-1-iux\,\mathbf{1}_{\{|x|\leq 1\}}
    \right)\nu(dx)
    \]
    encodes the entire jump structure of the process.
\end{itemize}

The L\'evy measure $\nu$ describes how frequently jumps of different sizes occur. The integrability condition \sourceref{eq:levy_measure_cond} has a precise meaning:
\begin{itemize}[label = --]
    \item on the region $\{|x|>1\}$ it implies
    \[
    \int_{\{|x|>1\}} \nu(dx)<\infty,
    \]
    so large jumps can occur only finitely many times on compact intervals;
    \item on the region $\{|x|\leq 1\}$ it imposes square-integrability of the small-jump contribution, which allows the compensated small-jump term to be well defined.
\end{itemize}

The truncation term
\[
-iux\,\mathbf{1}_{\{|x|\leq 1\}}
\]
is not an arbitrary correction. It is introduced to make the integral converge near the origin. Without this subtraction, the singular accumulation of small jumps would in general prevent the integral from being finite.


The general formula contains the principal examples already discussed.

\begin{itemize}[label = --]
    \item \textbf{Brownian motion with drift.} If $\nu=0$, the jump component is absent and
    \begin{equation}
    \psi(u)=i\gamma u-\frac{1}{2}Au^2,
    \end{equation}
    which is the characteristic exponent of the process $X_t = \gamma t + \sqrt{A}\,W_t$, consistently with~\sourceref{eq:BM_charfun}.
    %
    \item \textbf{Compound Poisson process.} If $A=0$ and
    \[
    \nu(dx)=\lambda F(dx),
    \]
    where $\lambda>0$ and $F$ is a probability measure on $\mathbb{R}\setminus\{0\}$, then the L\'evy measure has finite total mass $\nu(\mathbb{R}\setminus\{0\})=\lambda$, the truncation term can be absorbed into the drift, and one obtains
    \begin{equation}
    \psi(u)
    =
    i\gamma' u
    +
    \lambda
    \int_{\mathbb{R}\setminus\{0\}}
    \left(
    e^{iux}-1
    \right)F(dx),
    \end{equation}
    where $\gamma' = \gamma - \lambda\int_{\{|x|\leq 1\}} x\,F(dx)$ accounts for the recentring induced by the truncation convention. This is consistent with~\sourceref{eq:CPP_charfun}.
    %
    \item \textbf{Jump--diffusion.} If $A>0$ and $\nu(dx)=\lambda F(dx)$ with $\lambda>0$ and $F$ a probability measure on $\mathbb{R}\setminus\{0\}$, then the characteristic exponent combines the Gaussian and the compound Poisson components:
    \begin{equation}
    \psi(u)
    =
    i\gamma' u
    -\frac{1}{2}Au^2
    +
    \lambda
    \int_{\mathbb{R}\setminus\{0\}}
    \left(
    e^{iux}-1
    \right)F(dx).
    \end{equation}
    This is the L\'evy triplet that underlies the jump--diffusion dynamics developed above. The three terms reproduce, at the level of the characteristic exponent, the drift, the diffusive noise, and the compound Poisson jumps of the stochastic differential equation studied therein.
\end{itemize}

Thus Brownian motion, compound Poisson processes, and jump--diffusions are not external to the theory: they are embedded in the general L\'evy--Khintchine representation as particular specifications of the triplet $(\gamma,A,\nu)$.
\subsubsection{Why the formula matters in finance}

From the viewpoint of financial modelling, the L\'evy--Khintchine formula is important for two reasons.

First, it provides a complete classification of the admissible building blocks for return dynamics with stationary independent increments. Once the triplet
\[
(\gamma,A,\nu)
\]
is specified, the corresponding process is fully determined in law.

Second, it separates in a transparent way the three sources of variation:
\begin{itemize}[label = --]
    \item deterministic trend;
    \item continuous Gaussian fluctuation;
    \item discontinuous jump risk.
\end{itemize}

This decomposition is conceptually important because many asset-pricing models used in practice may be interpreted as different specifications of the L\'evy triplet. In particular, jump--diffusion models correspond to the case in which both $A\neq 0$ and the L\'evy measure $\nu$ is finite, whereas more general pure-jump models allow for infinite jump activity.

\begin{remark}
In asset-pricing applications the underlying price process is often modelled as $S_t = S_0 e^{X_t}$, where $\{X_t\}_{t\geq 0}$ is a L\'evy process representing the log-return. For the discounted price $e^{-rt}S_t$ to be a martingale under a risk-neutral measure, the exponent $\psi$ must satisfy
\[
\psi(-i) = r,
\]
which yields the risk-neutral drift restriction
\[
\gamma = r - \tfrac{1}{2}A - \int_{\mathbb{R}\setminus\{0\}}\left(e^x - 1 - x\,\mathbf{1}_{\{|x|\leq 1\}}\right)\nu(dx),
\]
provided $\int_{\{|x|>1\}} e^x\,\nu(dx)<\infty$. This condition generalises the drift compensation derived above for the jump--diffusion case.
\end{remark}

The L\'evy--Khintchine representation therefore closes the probabilistic block in a natural way: after introducing discontinuous dynamics through compound Poisson processes, we now arrive at the full class of continuous-time processes with stationary independent increments. In the following section, this structure becomes analytically even more useful, because characteristic functions and Fourier methods act directly on the exponent $\psi$, making valuation problems tractable even when transition densities are not available in closed form.

\section{Characteristic Functions and Fourier Methods}
\subsection{Characteristic Functions in Probability}
\subsubsection{The Intuitive and Analytical Framework}

In the classical Black-Scholes-Merton model, the probability density function (PDF) of the underlying asset's log-return is known in explicit closed form as a Gaussian distribution. This structural simplicity allows derivative payoffs to be priced by directly integrating over the real domain. However, as established in the preceding sections, transitioning to more realistic and complex market dynamics---such as jump-diffusions and general L\'evy processes---comes at a significant analytical cost: the exact probability density functions of these processes are rarely available in closed form.

To overcome this analytical bottleneck, modern quantitative finance shifts the pricing problem from the standard spatial domain (the state space of asset prices) into the frequency domain. This transformation relies heavily on the concept of the Characteristic Function (CF).

Conceptually, the characteristic function is the probabilistic analogue of the Fourier transform used in signal processing and systems engineering. Just as a physical signal can be decomposed into a continuous spectrum of complex exponential frequencies, a probability distribution can be completely identified by its Fourier transform. 

Analytically, for a random variable $X$ defined on a probability space $(\Omega, \mathcal{F}, \mathbb{P})$, the characteristic function $\varphi_X(u)$ is defined as the expected value of the complex exponential of $X$:
\begin{equation}
    \varphi_X(u) = \mathbb{E}[e^{iuX}] = \int_{-\infty}^{\infty} e^{iux} f_X(x) dx, \quad u \in \mathbb{R}
\end{equation}
where $i = \sqrt{-1}$ is the imaginary unit, $u$ is the real-valued frequency variable, and $f_X(x)$ represents the probability density function of $X$ (assuming it exists).

The fundamental mathematical advantage of the characteristic function over other transform methods, such as the Moment Generating Function (MGF), lies in the presence of the imaginary unit $i$. Because the magnitude of the complex exponential is always exactly one ($|e^{iux}| = 1$ for all real $u$ and $x$), the integral defining the characteristic function is absolutely convergent for \textit{any} valid probability distribution. This guarantees that the characteristic function always exists, even for heavy-tailed jump distributions (such as the Cauchy or Pareto distributions) where the mean, variance, or higher moments might violently diverge to infinity.

In the context of the L\'evy processes defined in Chapter 5, the characteristic function constitutes the ultimate analytical tool. While the transition density $f_{X_t}(x)$ of a L\'evy process $\{X_t\}_{t \ge 0}$ is generally unknown, its characteristic function is always explicitly known and structurally encoded by the characteristic exponent $\psi(u)$ via the L\'evy-Khintchine representation. Recalling the fundamental exponential representation from Chapter 5:
\begin{equation}
    \varphi_t(u) = \mathbb{E}[e^{iuX_t}] = e^{t\psi(u)}
\end{equation}

By exploiting Equation \sourceref{eq:levy_exponent_chap6}, complex convolution operations over disjoint time intervals mathematically collapse into simple algebraic multiplications in the frequency domain, rendering the valuation of contingent claims highly tractable.
\subsubsection{Moment Generation via the Characteristic Function}

Beyond ensuring analytical tractability in the frequency domain, the characteristic function serves as a direct mathematical mechanism for extracting the statistical moments of a random variable without requiring explicit knowledge of its density function. 

By expanding the complex exponential $e^{iuX}$ into its Taylor series representation and interchanging the expectation with the infinite sum (provided the necessary integrability conditions are met), one can establish a direct correspondence between the derivatives of the characteristic function evaluated at the origin ($u=0$) and the raw moments of the distribution.

Formally, if a random variable $X$ possesses finite absolute moments up to order $n$, that is $\mathbb{E}[|X|^n] < \infty$, then its characteristic function $\varphi_X(u)$ is $n$-times continuously differentiable on $\mathbb{R}$, and the $n$-th raw moment is given by:
\begin{equation}
    \mathbb{E}[X^n] = (-i)^n \left. \frac{d^n}{du^n} \varphi_X(u) \right|_{u=0}
\end{equation}

This property is profoundly advantageous when analyzing the continuous-time L\'evy processes $\{X_t\}_{t \ge 0}$ defined previously. Given the exponential representation in Equation \sourceref{eq:levy_exponent_chap6}, the moments of the asset return at any future time $t$ can be derived purely by differentiating the deterministic characteristic exponent $\psi(u)$. 

Applying the chain rule, the first derivative of the characteristic function with respect to $u$ is:
\begin{equation}
    \frac{d}{du} \varphi_t(u) = t \psi'(u) e^{t\psi(u)}
\end{equation}
Evaluating Equation \sourceref{eq:first_derivative_CF} at $u=0$ and noting that $\psi(0) = 0$ (which implies $e^{t\psi(0)} = 1$), the expected value (first moment) of the L\'evy process reduces to:
\begin{equation}
    \mathbb{E}[X_t] = -i \left. \frac{d}{du} \varphi_t(u) \right|_{u=0} = -i t \psi'(0)
\end{equation}

Similarly, computing the second derivative yields the second raw moment:
\begin{equation}
    \frac{d^2}{du^2} \varphi_t(u) = \left[ t \psi''(u) + \left( t \psi'(u) \right)^2 \right] e^{t\psi(u)}
\end{equation}
\begin{equation}
    \mathbb{E}[X_t^2] = (-i)^2 \left. \frac{d^2}{du^2} \varphi_t(u) \right|_{u=0} = - \left[ t \psi''(0) + t^2 (\psi'(0))^2 \right]
\end{equation}

From these raw moments, central moments such as the variance can be immediately constructed:
\begin{equation}
    Var(X_t) = \mathbb{E}[X_t^2] - (\mathbb{E}[X_t])^2 = -t \psi''(0)
\end{equation}

Consequently, measuring the expected drift, the variance, and even higher-order moments (such as skewness and kurtosis) of a complex jump-diffusion model does not require probability density integrations. The entire statistical profile of the market model can be parametrically controlled and extracted algebraically directly from the characteristic exponent $\psi(u)$.
\subsection{The Inversion Logic: Recovering Probabilities from Frequencies}

While the characteristic function guarantees mathematical tractability in the frequency domain, derivative pricing ultimately requires evaluating expected payoffs in the real spatial domain. This necessitates a robust mathematical procedure to "invert" the transformation and recover the probability distribution of the underlying asset.
\subsubsection{The Standard Fourier Inversion and its Computational Bottleneck}

If the characteristic function \(\varphi_X(u)\) of a random variable \(X\) is absolutely integrable, that is,
\[
\int_{-\infty}^{\infty} |\varphi_X(u)|\,du < \infty,
\]
then the probability distribution of \(X\) is absolutely continuous. Under this condition, the classical inverse Fourier transform yields the exact probability density function (PDF), \(f_X(x)\), evaluated at any real point \(x\):
\begin{equation}
    f_X(x) = \frac{1}{2\pi} \int_{-\infty}^{\infty} e^{-iux} \varphi_X(u) du
\end{equation}

While theoretically sound, relying on this standard inversion is computationally highly inefficient for derivative pricing. Pricing a European option natively requires the calculation of tail probabilities---specifically, the probability that the asset price exceeds the strike price at maturity, $\mathbb{P}(X > x)$. Utilizing the PDF would require a double numerical integration algorithm: an inner integral over the complex plane to invert the Fourier transform, followed by an outer integral over the real domain to compute the Cumulative Distribution Function (CDF). In calibration routines where prices must be evaluated thousands of times, this bottleneck becomes computationally prohibitive.
\subsubsection{The Gil-Pelaez Inversion Theorem}

To circumvent the computational burden of double integration, quantitative finance heavily relies on the Gil-Pelaez Theorem (1951). This powerful analytical result establishes a direct relationship between the characteristic function and the Cumulative Distribution Function, $F_X(x) = \mathbb{P}(X \le x)$, requiring only a single, well-behaved integration over the positive real line.

For any continuous point $x$ of the distribution, the Gil-Pelaez inversion formula states:
\begin{equation}
    F_X(x) = \mathbb{P}(X \le x) = \frac{1}{2} - \frac{1}{\pi} \int_{0}^{\infty} \Re \left[ \frac{e^{-iux} \varphi_X(u)}{iu} \right] du
\end{equation}
where $\Re[\cdot]$ denotes the real part of the complex argument.

Consequently, the complementary tail probability---which forms the fundamental building block for European call options---is immediately derived as:
\begin{equation}
    \mathbb{P}(X > x) = 1 - F_X(x) = \frac{1}{2} + \frac{1}{\pi} \int_{0}^{\infty} \Re \left[ \frac{e^{-iux} \varphi_X(u)}{iu} \right] du
\end{equation}
\subsubsection{Application to Risk-Neutral Exercise Probabilities}

In the context of the L\'evy-driven asset dynamics discussed in Chapter 5, let $X_t = \ln(S_t)$ be the logarithmic price process. Under the risk-neutral measure, the Gil-Pelaez theorem allows for the direct computation of the exact exercise probability of a Call option maturing at time $t$ with a logarithmic strike price $k = \ln(K)$.

By substituting the exponential representation of the L\'evy process characteristic function from Equation \sourceref{eq:levy_exponent_chap6} into the inversion formula, the risk-neutral probability of the option expiring in-the-money is exactly:
\begin{equation}
    \mathbb{P}(\ln S_t > k) = \frac{1}{2} + \frac{1}{\pi} \int_{0}^{\infty} \Re \left[ \frac{e^{-iuk} e^{t\psi(u)}}{iu} \right] du
\end{equation}

This single-integral formulation is the cornerstone of modern semi-analytical pricing methods. By reducing the dimensionality of the integration, it ensures that models with arbitrary jump structures and stochastic volatility can be evaluated and calibrated to market data with computational speeds comparable to the closed-form Black-Scholes formulas.
\subsection{Fourier Transforms in Pricing Theory}

The Gil-Pelaez inversion discussed in Section 6.2 is highly effective for pricing binary (digital) options, as their payoffs mathematically correspond to simple tail probabilities. However, pricing standard "vanilla" derivatives, such as European Call options, introduces a significant analytical obstacle related to the strict integrability requirements of Fourier analysis.
\subsubsection{The Integrability Problem of Option Payoffs}

Let $S_T$ be the asset price at maturity $T$, and let $s_T = \ln(S_T)$ be the terminal logarithmic price. Consider a European Call option with strike $K$, and define the logarithmic strike as $k = \ln(K)$. Under the risk-neutral measure $\mathbb{Q}$, the present value of the Call option, denoted as $C(k)$, is the discounted expected value of its terminal payoff:
\begin{equation}
    C(k) = e^{-rT} \mathbb{E}^{\mathbb{Q}}[\max(S_T - K, 0)] = e^{-rT} \int_{k}^{\infty} (e^{s_T} - e^k) q_T(s_T) ds_T
\end{equation}
where $r$ is the constant risk-free rate, and $q_T(s_T)$ is the risk-neutral probability density function of the terminal log-price.

To leverage the speed and tractability of the frequency domain, one would intuitively attempt to compute the analytical Fourier transform of the pricing function $C(k)$ with respect to the log-strike $k$. The Fourier transform of $C(k)$ would be defined as:
\begin{equation}
    \hat{C}(u) = \int_{-\infty}^{\infty} e^{iuk} C(k) dk
\end{equation}

For this integral to exist and be well-defined, the function $C(k)$ must be absolutely integrable, meaning $C(k) \in \mathcal{L}^1(\mathbb{R})$, such that $\int_{-\infty}^{\infty} |C(k)| dk < \infty$. 

However, evaluating the asymptotic behavior of the Call price function violently violates this condition. As the log-strike approaches positive infinity ($k \to \infty$), the option becomes deeply out-of-the-money, and the price correctly decays to zero ($C(k) \to 0$). Conversely, as the log-strike approaches negative infinity ($k \to -\infty$), the option becomes deeply in-the-money. In this limit, the strike price $K$ tends to zero, and the value of the Call option approaches the initial asset price:
\begin{equation}
    \lim_{k \to -\infty} C(k) = S_0
\end{equation}

Because the function converges to a strictly positive constant rather than zero as $k \to -\infty$, the area under the curve $C(k)$ is infinite. Therefore, $C(k) \notin \mathcal{L}^1(\mathbb{R})$, and its standard Fourier transform $\hat{C}(u)$ simply does not exist. This lack of integrability prevents the direct application of Fourier pricing techniques to vanilla options.
\subsubsection{The Dampening Factor (The Carr-Madan Approach)}

To resolve the non-integrability of the Call option price, Carr and Madan (1999) introduced a mathematically elegant modification: multiplying the option price by an exponential dampening factor. This technique conceptually mirrors the transition from the Fourier transform to the Laplace transform in signal processing, where a divergent signal is forced to decay exponentially.

Let $\alpha > 0$ be a strictly positive damping coefficient. We define the dampened, or modified, Call option price $c(k)$ as:
\begin{equation}
    c(k) = e^{\alpha k} C(k)
\end{equation}



The introduction of the exponential term $e^{\alpha k}$ perfectly stabilizes the asymptotic behavior of the function on both ends of the real line:
\begin{itemize}[label = --]
    \item \textbf{As $k \to -\infty$ (Deep In-The-Money):} The Call price $C(k)$ is bounded by the initial asset price $S_0$. The damping factor $e^{\alpha k}$ rapidly decays to zero. Consequently, the modified price vanishes: $\lim_{k \to -\infty} c(k) = 0$.
    \item \textbf{As $k \to \infty$ (Deep Out-Of-The-Money):} The damping factor $e^{\alpha k}$ grows exponentially. However, for a properly chosen $\alpha$, the probability density of the terminal log-price decays faster than $e^{\alpha k}$ grows, forcing the product to zero: $\lim_{k \to \infty} c(k) = 0$.
\end{itemize}

Because the dampened pricing function $c(k)$ approaches zero at both extremities, it becomes absolutely integrable, i.e., $c(k) \in \mathcal{L}^1(\mathbb{R})$. 

% Asset/code block omitted pending provenance acceptance.


With integrability successfully restored, we can now compute the standard Fourier transform of the dampened price $c(k)$ with respect to the log-strike $k$. Let $\psi_T(u)$ denote this Fourier transform:
\begin{equation}
    \psi_T(u) = \int_{-\infty}^{\infty} e^{iuk} c(k) dk = \int_{-\infty}^{\infty} e^{iuk} e^{\alpha k} C(k) dk
\end{equation}
which can be compactly rewritten by combining the exponents:
\begin{equation}
    \psi_T(u) = \int_{-\infty}^{\infty} e^{i(u - i\alpha)k} C(k) dk
\end{equation}

This equation demonstrates that evaluating the standard Fourier transform of the dampened option price along the real axis $u$ is mathematically equivalent to evaluating the Fourier transform of the original undampened option price along a shifted contour in the complex plane, specifically at the complex frequency $(u - i\alpha)$.
\subsubsection{The Analytical Pricing Formula}

The true power of the Carr-Madan dampening technique lies in the explicit connection it establishes between the Fourier transform of the option price and the characteristic function of the underlying asset's return. 

By substituting the risk-neutral pricing expectation into the definition of the dampened transform $\psi_T(u)$ and evaluating the spatial integrals, Carr and Madan proved that the transform of the dampened Call price can be expressed entirely as an algebraic function of the terminal log-price characteristic function. The exact analytical relation is:
\begin{equation}
    \psi_T(u) = \frac{e^{-rT} \varphi_T(u - i(\alpha + 1))}{\alpha^2 + \alpha - u^2 + i(2\alpha + 1)u}
\end{equation}

Here, $\varphi_T(z) = \mathbb{E}^{\mathbb{Q}}[e^{iz s_T}]$ is the characteristic function of the terminal logarithmic price $s_T = \ln(S_T)$ evaluated at the complex argument $z = u - i(\alpha + 1)$. For any underlying L\'evy process, this characteristic function is inherently known in closed form via the L\'evy-Khintchine exponent, ensuring that the numerator $\varphi_T(\cdot)$ is readily computable.

Once the analytical form of $\psi_T(u)$ is established in the frequency domain, the final step is to recover the undampened, real-world Call option price $C(k)$. This is achieved by applying the Inverse Fourier Transform to $\psi_T(u)$ to obtain the dampened price $c(k)$, and subsequently multiplying by the inverse dampening factor $e^{-\alpha k}$. 

Since the Call option price must be a strictly real number, the inverse Fourier integral can be significantly simplified. The imaginary components must mathematically cancel out, allowing the integration to be performed exclusively over the positive real half-line by taking the real part of the integrand:
\begin{equation}
    C(k) = \frac{e^{-\alpha k}}{\pi} \int_{0}^{\infty} \operatorname{Re}\!\left[ e^{-iuk}\psi_T(u) \right]du.
\end{equation}

This final equation represents a paradigm shift in quantitative finance. By substituting the known algebraic expression for $\psi_T(u)$ into the integral, the pricing of a European Call option under highly complex, discontinuous, or stochastic volatility dynamics is reduced to the numerical evaluation of a single, highly stable, one-dimensional integral.

% ---- next approved source slice ----

\section{Stochastic Volatility Theory}

The structural failures of the constant-volatility hypothesis have already
been recorded at the beginning of this chapter: the implied volatility surface is not flat,
log-returns are leptokurtic, and the realised quadratic variation of the
underlying does not behave as a deterministic function of time. The present
chapter develops the analytical framework that removes the rigidity of a
constant diffusion coefficient by promoting the instantaneous variance to
a non-anticipative stochastic process in its own right. The construction
is carried out first in full generality, in order to clarify the
probabilistic content of the extension, and is subsequently specialised
to the canonical example of Heston, whose tractability rests on a
square-root specification analysed below within the broader theory
of affine diffusions.
\subsection{Empirical Background and Two-Factor Structure}


The empirical evidence that motivates the abandonment of a constant
diffusion coefficient can be summarised in four stylised facts, none of
which is reproducible inside the Black--Scholes--Merton framework
\sourceref{eq:BS_SDE}.

\medskip

\textit{- Strike-dependent implied volatilities.} Section~2.2 already
established that for a fixed maturity $\tau$ the implied volatility
$\sigma_{\mathrm{imp}}(K,\tau)$ extracted from quoted European options is
not a constant function of the strike $K$. In equity index markets the
function $K\mapsto \sigma_{\mathrm{imp}}(K,\tau)$ is monotonically
decreasing and convex, while in foreign exchange markets it displays a
symmetric U-shaped profile. In both cases the surface is incompatible
with any model in which the diffusion coefficient is a deterministic
constant.

\textit{- Maturity-dependent shape.} The skew is steepest for short
maturities and progressively flattens as $\tau$ grows. Any single
parameter $\sigma>0$ is therefore unable to reproduce the joint dependence
of $\sigma_{\mathrm{imp}}$ on $(K,\tau)$.

\textit{- Volatility clustering.} Realised variance computed on rolling
windows of intra-day returns exhibits significant positive autocorrelation
over horizons ranging from days to several months: large absolute returns
are statistically followed by large absolute returns, and quiet periods
are followed by quiet periods. This time-series structure is excluded by
the geometric Brownian dynamics, under which squared log-returns are
independent.

\textit{- Leverage effect.} Innovations in price and innovations in
volatility are negatively correlated in equity markets: a downward
movement of the underlying is typically accompanied by an upward revision
of the variance, while upward movements are followed by smaller, often
negligible, downward revisions of volatility. Such an asymmetric
dependence cannot exist within a one-dimensional diffusion driven by a
single Brownian motion.

\medskip

The minimal modification of the Black--Scholes dynamics consistent with
(a)--(d) consists in enlarging the state space by adjoining a second
factor that represents the instantaneous variance and that is itself
governed by a stochastic differential equation, possibly correlated with
the Brownian motion driving the price. To formalise this idea we work on
the filtered probability space
$(\Omega,\mathcal{F},\{\mathcal{F}_t\}_{t\geq 0},\mathbb{P})$ already
introduced in the preceding sections, satisfying the usual conditions of
right-continuity and $\mathbb{P}$-completeness. On this space we consider
two scalar standard Brownian motions $\{W^{S}_t\}_{t\geq 0}$ and
$\{W^{v}_t\}_{t\geq 0}$, both adapted to $\{\mathcal{F}_t\}_{t\geq 0}$,
with constant instantaneous correlation
\begin{equation}
d\langle W^{S},W^{v}\rangle_t \;=\; \rho\,dt,
\qquad \rho\in[-1,1].
\end{equation}
Equivalently, one may write
\begin{equation}
W^{v}_t
=
\rho\,W^{S}_t+\sqrt{1-\rho^{2}}\,Z_t,
\end{equation}
where $\{Z_t\}_{t\geq 0}$ is a standard Brownian motion independent of
$\{W^{S}_t\}_{t\geq 0}$.

\begin{definition}[Stochastic volatility model]
A two-factor diffusion model on
$(\Omega,\mathcal{F},\{\mathcal{F}_t\}_{t\geq 0},\mathbb{P})$ is called a
\emph{stochastic volatility model} for the asset price
$\{S_t\}_{t\geq 0}$ if there exist measurable coefficients $\mu$,
$\alpha$, $\beta$ and a strictly positive,
$\{\mathcal{F}_t\}_{t\geq 0}$-progressively measurable process
$\{v_t\}_{t\geq 0}$ such that
\begin{align}
dS_t &= \mu(t,S_t,v_t)\,S_t\,dt
       + \sqrt{v_t}\,S_t\,dW^{S}_t,
       \\
dv_t &= \alpha(t,v_t)\,dt
       + \beta(t,v_t)\,dW^{v}_t,
\end{align}
with $S_0>0$, $v_0>0$ and $W^{S}, W^{v}$ correlated as in
\sourceref{eq:SV_corr}, and provided that the system
\sourceref{eq:SV_price}--\sourceref{eq:SV_variance} admits a unique strong
solution on $[0,T]$ for every finite horizon $T>0$.
\end{definition}

The state vector of the model is the pair $(S_t,v_t)$. The price equation
\sourceref{eq:SV_price} replaces the constant diffusion coefficient $\sigma$
of \sourceref{eq:BS_SDE} with the random non-anticipative volatility
$\sqrt{v_t}$, while \sourceref{eq:SV_variance} specifies the dynamics of the
variance and characterises the particular member of the class. Different
choices of $(\alpha,\beta)$ produce qualitatively different models: a
linear drift $\alpha(t,v) = \kappa(\theta-v)$ with constants
$\kappa,\theta>0$ enforces mean reversion of the variance to the long-run
level $\theta$ at rate $\kappa$, while the diffusion coefficient $\beta$
controls the magnitude of the random fluctuations of $v_t$ around its
drift.

\begin{remark}
The class of stochastic volatility models is structurally distinct from
the class of \emph{local-volatility} models, in which the diffusion of
the price is a deterministic function of $(t,S_t)$. Local volatility
introduces no additional source of randomness beyond the Brownian motion
already driving the price; it can fit a given implied volatility surface
at a fixed observation date, but it does not generate an autonomous
time-series dynamics for variance. Stochastic volatility, on the
contrary, is driven by the genuinely two-dimensional Brownian noise
$(W^{S},W^{v})$: the variance is itself random, the source of risk is
intrinsically two-dimensional, and the analytical price to pay for this
enrichment is the loss of completeness, which we examine in the next
section.
\end{remark}
\subsection{Risk-Neutral Dynamics, Correlation and the Leverage Effect}


Within \sourceref{eq:SV_price}--\sourceref{eq:SV_variance} the asset $S_t$ is a
tradable quantity, while the variance $v_t$ is not directly traded on
any market. The two-dimensional Brownian noise driving the system
therefore carries two sources of risk, but only one liquid traded asset
is available for dynamic hedging. This asymmetry is the structural origin
of market incompleteness in stochastic volatility models, in direct
analogy with the incompleteness already encountered in Section~4.4 for
jump--diffusions, and it determines both the construction of the pricing
measure and the economic interpretation of the additional state variable.

To make the statement precise, it is convenient to work with the
independent Brownian representation \sourceref{eq:SV_BM_decomposition}. The
independent sources of randomness are then $W^{S}$ and $Z$, while the
variance Brownian motion is recovered from
\[
W^{v}_t
=
\rho W^{S}_t+\sqrt{1-\rho^2}Z_t.
\]
Let $r$ denote the constant risk-free rate and assume an equivalent
change of measure $\mathbb{P}\sim\mathbb{Q}$ characterised by the
Girsanov density
\begin{equation}
\frac{d\mathbb{Q}}{d\mathbb{P}}\bigg|_{\mathcal{F}_t}
=
\exp\!\left\{
-\int_0^t \lambda^S_s\,dW^S_s
-\int_0^t \lambda^Z_s\,dZ_s
-\frac12\int_0^t
\left((\lambda^S_s)^2+(\lambda^Z_s)^2\right)\,ds
\right\},
\end{equation}
where $\lambda^S$ and $\lambda^Z$ are progressively measurable processes
satisfying a suitable integrability condition, for instance Novikov's
condition. Under $\mathbb{Q}$ the processes
\begin{equation}
dW^{S,\mathbb{Q}}_t
=
dW^{S}_t+\lambda^S_t\,dt,
\qquad
dZ^{\mathbb{Q}}_t
=
dZ_t+\lambda^Z_t\,dt
\end{equation}
are independent standard Brownian motions. The corresponding
risk-neutral Brownian motion driving the variance is
\begin{equation}
dW^{v,\mathbb{Q}}_t
=
\rho\,dW^{S,\mathbb{Q}}_t
+
\sqrt{1-\rho^2}\,dZ^{\mathbb{Q}}_t,
\end{equation}
so that the instantaneous correlation between $W^{S,\mathbb{Q}}$ and
$W^{v,\mathbb{Q}}$ remains equal to $\rho$.

The discounted price $e^{-rt}S_t$ is a $\mathbb{Q}$-local martingale if
and only if the drift of the stock price under $\mathbb{Q}$ is equal to
$r$. Since the stock price equation is driven directly by $W^{S}$, this
condition fixes the market price of risk associated with the tradable
source of randomness:
\begin{equation}
\mu(t,S_t,v_t)-r
=
\sqrt{v_t}\,\lambda^S_t.
\end{equation}
Equation \sourceref{eq:SV_drift_restriction} pins down the component
$\lambda^S$, the classical market price of price risk: it is determined
by no-arbitrage once the physical drift $\mu$ is given.

The second component $\lambda^Z$ is not pinned down by no-arbitrage,
because it is associated with the source of randomness that affects the
variance but is not directly spanned by trading in the underlying asset
alone. Equivalently, after rewriting the system in terms of the
correlated Brownian pair $(W^{S,\mathbb{Q}},W^{v,\mathbb{Q}})$, there
remains one free degree of freedom in the drift of the variance process
under $\mathbb{Q}$. This free component is called the \emph{market price
of volatility risk}. In practice it is selected so that the variance
process retains a tractable parametric form under $\mathbb{Q}$, typically
of the same family as under $\mathbb{P}$. Once this choice is made, the
model becomes \emph{calibrated} rather than purely \emph{estimated}: the
risk-neutral parameters are extracted from the cross-section of option
prices and absorb the volatility risk premium required by the market.

Throughout the remainder of the chapter we shall work directly under a
chosen pricing measure $\mathbb{Q}$ and denote the corresponding
correlated Brownian motions simply by $W^{S}$ and $W^{v}$, with the
understanding that all coefficients are risk-neutral ones. The
risk-neutral price equation is then
\begin{equation}
dS_t
=
r\,S_t\,dt+\sqrt{v_t}\,S_t\,dW^{S}_t.
\end{equation}

Once the pricing measure is fixed, the geometry of the implied volatility
surface generated by \sourceref{eq:SV_price_RN} and
\sourceref{eq:SV_variance} is governed in first instance by the correlation
parameter $\rho$ in \sourceref{eq:SV_corr}. This parameter is the channel
through which innovations in the variance feed back into innovations in
the price. Its sign and magnitude have a direct effect on the asymmetry
of the terminal log-return distribution and, consequently, on the shape
of the implied volatility surface.

\begin{definition}[Uncorrelated variance shocks]
Suppose that $\rho=0$. Then the Brownian motion driving the stock price
is independent of the Brownian motion driving the variance. Conditionally
on the path $\{v_s\}_{s\in[0,T]}$, the terminal log-price
$\log(S_T/S_0)$ is normally distributed with mean
$rT-\frac12 I_T$ and variance $I_T$, where
\begin{equation}
I_T
:=
\int_0^T v_s\,ds.
\end{equation}
Equivalently,
\begin{equation}
\log(S_T/S_0)\,\big|\,\mathcal{F}^{v}_{T}
\sim
\mathcal{N}\!\left(
rT-\frac12 I_T,\,
I_T
\right),
\qquad
\mathcal{F}^{v}_{T}:=\sigma(v_s:0\leq s\leq T).
\end{equation}
\end{definition}

\begin{proof}[Sketch]
Applying the It\^o--Doeblin formula to
\sourceref{eq:SV_price_RN}, with $X_t:=\log S_t$, gives
\begin{equation}
dX_t
=
\left(r-\frac12 v_t\right)dt+\sqrt{v_t}\,dW^S_t.
\end{equation}
Integrating over $[0,T]$ yields
\begin{equation}
X_T-X_0
=
rT-\frac12\int_0^T v_s\,ds
+
\int_0^T \sqrt{v_s}\,dW^S_s.
\end{equation}
When $\rho=0$, the Brownian motion $W^S$ is independent of the variance
filtration $\mathcal{F}^{v}_{T}$. Therefore, conditionally on the path of
$v$, the stochastic integral in \sourceref{eq:SV_log_integrated} is Gaussian
with mean zero and variance $\int_0^T v_s\,ds$. This gives
\sourceref{eq:SV_conditional_logreturn}.
\end{proof}

The unconditional law of the log-return is therefore a Gaussian mixture
indexed by the random integrated variance $I_T$. This mechanism produces
heavier tails than the Black--Scholes Gaussian benchmark, because large
values of $I_T$ generate wider conditional Gaussian distributions. Exact
symmetry of the implied volatility smile, however, does not follow in
full generality from $\rho=0$, since the conditional mean in
\sourceref{eq:SV_conditional_logreturn} also depends on $I_T$ through the
term $-\frac12 I_T$. The robust conclusion is instead that uncorrelated
stochastic volatility mainly produces excess kurtosis and smile
curvature, while the leverage-induced negative skew observed in equity
index options requires a non-zero, and typically negative, correlation.

\begin{remark}
The terminology \emph{leverage effect} is historically due to the
heuristic argument that a fall in the equity value of a firm increases
the debt-to-equity ratio and therefore the financial risk borne by the
remaining shareholders, with a corresponding increase in equity
volatility. Although the empirical magnitude of the negative correlation
between equity returns and volatility is too large to be fully explained
by changes in capital structure alone, the term is universally used to
denote the phenomenon itself. Within
\sourceref{eq:SV_corr}--\sourceref{eq:SV_price_RN}, the leverage effect is
encoded by $\rho<0$: negative shocks to the stock price tend to coincide
with positive shocks to the variance. This negative correlation is the
principal mechanism through which stochastic volatility models generate
the downward implied-volatility skew observed in equity index markets.
\end{remark}

A second economic consequence of market incompleteness deserves to be
mentioned at this stage. The free component in the change of measure
\sourceref{eq:SV_girsanov} represents the compensation demanded by market
participants for bearing innovations in $v_t$ that cannot be perfectly
hedged by trading only in the underlying. This compensation is usually
called the \emph{variance risk premium}. In equity markets the variance
risk premium is typically found to be negative in empirical studies:
investors are willing to pay a premium for protection against upward
variance shocks. This observation is consistent with the distinction
between physical and risk-neutral variance dynamics.
\subsection{Vol-of-vol and the Heston Specification}


The correlation parameter determines the direction of the leverage
feedback, but its magnitude must be combined with that of the diffusion
coefficient $\beta$ in \sourceref{eq:SV_variance} in order to control the
overall steepness and curvature of the implied volatility surface. The
coefficient $\beta$ is referred to as the \emph{volatility of
volatility}, or \emph{vol-of-vol}, since it measures the intensity of the
random oscillations of the variance around its drift. It acts on the
smile through two distinct channels.

\begin{itemize}[label = --]
    \item \textit{Smile curvature.} Conditionally on the trajectory of
          the variance, the terminal log-price is Gaussian, but its
          unconditional distribution is a mixture of Gaussians indexed
          by the random integrated variance $I_T$ defined in
          \sourceref{eq:integrated_variance}. This mechanism generates
          heavier tails than the Black--Scholes lognormal benchmark.

          To see the mechanism in a transparent way, consider the
          simplified centred mixture
          \[
          Y_T=\sqrt{I_T}\,Z,
          \qquad Z\sim\mathcal{N}(0,1),
          \]
          with $Z$ independent of $I_T$. In this simplified setting one
          obtains
          \begin{equation}
          \operatorname{Kurt}(Y_T)
          =
          3+
          \frac{3\,\mathrm{Var}(I_T)}
          {\bigl(\mathbb{E}[I_T]\bigr)^2}.
          \end{equation}
          Formula \sourceref{eq:SV_kurtosis} should be interpreted as a
          structural calculation rather than as the exact kurtosis
          formula for $\log S_T$, because the true log-price contains
          the additional drift correction $-\frac12 I_T$. Nevertheless,
          it captures the essential point: randomness of integrated
          variance produces excess kurtosis relative to the
          Black--Scholes benchmark. Increasing the vol-of-vol raises the
          variability of $I_T$, fattens the tails of the return
          distribution, and deepens the implied volatility smile.

    \item \textit{Skew amplification.} The vol-of-vol amplifies the
          asymmetry generated by a non-zero $\rho$. The correlation
          parameter determines the sign of the price--variance feedback,
          while the vol-of-vol determines the size of the variance
          shocks that can be correlated with price shocks. Thus a
          negative $\rho$ generates the direction of the equity skew,
          whereas a large vol-of-vol increases its magnitude. If the
          vol-of-vol is very small, even a strongly negative correlation
          has limited effect on the implied volatility surface, because
          the variance process itself moves very little.
\end{itemize}

When the vol-of-vol vanishes, the variance process collapses toward its
deterministic mean-reverting limit, the Gaussian-mixture effect
disappears, and the implied volatility surface moves back toward the flat
Black--Scholes surface. The structural role of the variance process can
therefore be summarised as follows: stochastic volatility generates
excess kurtosis because the integrated variance $I_T$ is random rather
than deterministic, and it generates skew through the interaction between
variance shocks and price shocks.

Among the parametric specifications of
\sourceref{eq:SV_price}--\sourceref{eq:SV_variance}, the most influential one is
obtained by combining a linear mean-reverting drift with a square-root
diffusion coefficient for the variance. This choice is particularly
important because it preserves non-negativity, generates leverage and
excess kurtosis, and retains the affine structure needed for
transform-based pricing.

\begin{definition}[Heston model]
The \emph{Heston stochastic volatility model} is the two-factor diffusion
defined under the risk-neutral measure $\mathbb{Q}$ by
\begin{equation}
dS_t
=
r\,S_t\,dt+\sqrt{v_t}\,S_t\,dW^{S}_t,
\end{equation}
\begin{equation}
dv_t
=
\kappa(\theta-v_t)\,dt+\xi\sqrt{v_t}\,dW^{v}_t,
\end{equation}
\begin{equation}
d\langle W^{S},W^{v}\rangle_t
=
\rho\,dt,
\end{equation}
with $S_0>0$, $v_0>0$, and constants $\kappa>0$, $\theta>0$, $\xi>0$,
$\rho\in[-1,1]$.
\end{definition}

The four parameters of \sourceref{eq:heston_S}--\sourceref{eq:heston_corr} admit
a transparent interpretation. Taking expectations on both sides of
\sourceref{eq:heston_v} and exploiting the martingale property of the
stochastic integral yields
\[
\frac{d}{dt}\mathbb{E}^{\mathbb{Q}}[v_t]
=
\kappa\left(\theta-\mathbb{E}^{\mathbb{Q}}[v_t]\right).
\]
Solving this ordinary differential equation gives
\begin{equation}
\mathbb{E}^{\mathbb{Q}}[v_t]
=
\theta+(v_0-\theta)e^{-\kappa t}.
\end{equation}
The expected variance therefore converges exponentially to $\theta$,
which has the interpretation of the long-run level of instantaneous
variance. The parameter $\kappa$ controls the rate of convergence in
\sourceref{eq:heston_meanvar}: large values of $\kappa$ correspond to a
variance process that adjusts quickly to its long-run level, while small
values correspond to a persistent process whose deviations from $\theta$
decay slowly. The half-life of a variance shock is $\log 2/\kappa$.

The parameter $\xi$, which plays the role of the vol-of-vol introduced
above, makes the instantaneous variance of changes in $v_t$ proportional
to $\xi^2v_t$: shocks to the variance are larger when the variance itself
is high. Finally, the parameter $\rho$ controls the instantaneous
dependence between shocks to the price and shocks to the variance.
Negative values of $\rho$ encode the leverage effect and are the
principal source of the negative implied-volatility skew observed in
equity index markets.

The square-root structure of \sourceref{eq:heston_v} is introduced because it
combines economic interpretability and analytical tractability. The drift
term $\kappa(\theta-v_t)$ produces mean reversion, while the diffusion
coefficient $\xi\sqrt{v_t}$ vanishes at the boundary $v=0$, helping to
preserve non-negativity of the variance process. The precise boundary
classification, the existence and uniqueness of the strong solution, and
the role of the Feller condition
\[
2\kappa\theta\geq \xi^2
\]
are deferred to the affine comparison below, where \sourceref{eq:heston_v} is studied as the
Cox--Ingersoll--Ross square-root diffusion.

\begin{remark}
Other parametric stochastic volatility specifications appear in the
literature. The lognormal specification of Hull and White, the
Stein--Stein model based on an Ornstein--Uhlenbeck process for
volatility, the inverse-square-root or $3/2$ model, and the SABR model
are all instances of the general two-factor framework
\sourceref{eq:SV_price}--\sourceref{eq:SV_variance}, corresponding to different
choices of the variance or volatility dynamics. The Heston model is
especially important because it combines a financially interpretable
variance process with the affine structure needed for transform-based
option pricing.
\end{remark}
\subsection{Pricing via the Characteristic Function}


The decisive analytical advantage of the Heston specification is that
the characteristic function of the terminal log-price is available in
closed form, which reduces option pricing to the transform machinery
already developed earlier in this chapter. Set $X_t:=\log S_t$ and let
\begin{equation}
\varphi_T(u)
:=
\mathbb{E}^{\mathbb{Q}}
\left[
e^{iuX_T}\mid\mathcal{F}_0
\right],
\qquad u\in\mathbb{R},
\end{equation}
denote the characteristic function of $X_T$ under $\mathbb{Q}$. For the
dynamics \sourceref{eq:heston_S}--\sourceref{eq:heston_corr}, $\varphi_T(u)$
admits the exponential-affine representation
\begin{equation}
\varphi_T(u)
=
\exp\left\{
A(T,u)+B(T,u)v_0+iuX_0
\right\},
\end{equation}
where the deterministic functions $A(\cdot,u)$ and $B(\cdot,u)$ solve a
system of ordinary differential equations of Riccati type whose explicit
solution is provided above. The structural form
\sourceref{eq:heston_CF_affine} is a consequence of the affine dependence of
the drift and covariance structure of the state vector $(X_t,v_t)$ on
the state variables themselves. It is the analytical fingerprint of the
broader class of affine processes to which the Heston model belongs.

The availability of \sourceref{eq:heston_CF_affine} reduces the pricing of
European options under \sourceref{eq:heston_S}--\sourceref{eq:heston_corr} to the
framework developed earlier in this chapter. Combining
\sourceref{eq:heston_CF_affine} with the Carr--Madan transform
\sourceref{eq:carr_madan_pricing}, the price at $t=0$ of a European call
with strike $K=e^{k}$ and maturity $T$ is
\begin{equation}
C(k)
=
\frac{e^{-\alpha k}}{\pi}
\int_0^{\infty}
\Re\!\left[
e^{-iuk}
\frac{
e^{-rT}\,
\varphi_T\!\left(u-i(\alpha+1)\right)
}{
\alpha^{2}+\alpha-u^{2}+i(2\alpha+1)u
}
\right]du,
\end{equation}
with $\alpha>0$ a damping coefficient chosen in such a way that
$\varphi_T(u-i(\alpha+1))$ is well defined. Equivalently, the exercise
probabilities required for the writing of the price in the canonical
Black--Scholes form
\[
C=S_0P_1-Ke^{-rT}P_2
\]
can be computed via the Gil--Pelaez inversion
\sourceref{eq:gil_pelaez_tail} applied separately to the log-price and to
the share-measure characteristic functions, both of which are explicitly
known via \sourceref{eq:heston_CF_affine}.

\begin{remark}
The fact that \sourceref{eq:heston_call_price} reduces the pricing problem
to the numerical evaluation of a single one-dimensional integral is the
operational reason for the widespread use of the Heston model in
practice. Calibration to a cross-section of liquid option quotes, which
in a generic stochastic volatility model would require Monte Carlo
simulation of \sourceref{eq:SV_price}--\sourceref{eq:SV_variance} for each
parameter configuration, is reduced under
\sourceref{eq:heston_S}--\sourceref{eq:heston_corr} to the repeated numerical
computation of \sourceref{eq:heston_call_price}. The model is therefore rich
enough to generate smile and skew effects, but still tractable enough
for calibration to option surfaces.
\end{remark}
\subsection{Limits of Pure Stochastic Volatility}


The stochastic-volatility framework developed in this chapter removes the
constant-volatility restriction of Black--Scholes by replacing the fixed
parameter $\sigma$ with the random variance process $v_t$.  This
extension is sufficient to generate volatility clustering, excess
kurtosis, leverage effects and non-flat implied volatility surfaces.
Nevertheless, models of the form
\sourceref{eq:SV_price}--\sourceref{eq:SV_variance} remain purely diffusive:
their sample paths are almost surely continuous and therefore contain no
jump component.  In particular,
\begin{equation}
\Delta S_t:=S_t-S_{t-}=0
\qquad\text{almost surely}.
\end{equation}

This pathwise property becomes especially restrictive at short
maturities.  Under the risk-neutral dynamics
\sourceref{eq:SV_price_RN}, the log-price satisfies
\[
d\log S_t
=
\left(r-\frac12 v_t\right)dt+\sqrt{v_t}\,dW^S_t.
\]
Hence
\begin{equation}
\log\frac{S_T}{S_0}
=
rT-\frac12 I_T
+
\int_0^T \sqrt{v_s}\,dW^S_s,
\qquad
I_T:=\int_0^T v_s\,ds.
\end{equation}
Conditionally on the variance path, the stochastic integral in
\sourceref{eq:SV_log_short} is Gaussian with variance $I_T$.  If the
variance process is regular at the origin, then
\[
I_T=v_0T+o(T),
\qquad T\downarrow 0,
\]
and consequently
\begin{equation}
\log\frac{S_T}{S_0}
=
\left(r-\frac12 v_0\right)T
+
\sqrt{v_0}\,W_T
+
o_{\mathbb{P}}(\sqrt{T}).
\end{equation}
Thus, at leading order, the short-time behaviour of a pure stochastic
volatility diffusion is still locally Gaussian.  The randomness of
$v_t$ affects higher-order corrections, but it does not create an
order-one discontinuous return over an arbitrarily short horizon.

\begin{definition}[Short-maturity limitation of pure stochastic volatility]

A pure diffusion-based stochastic volatility model is said to exhibit a
\emph{short-maturity limitation} when its continuous-path structure makes
the leading short-time behaviour of log-returns locally Gaussian, so that
very steep short-dated smiles and skews can be reproduced only through
higher-order variance effects or through extreme parameter values.
\end{definition}

This limitation is structural, not specific to the Heston
parametrisation.  The Heston model enriches Black--Scholes by making the
variance stochastic and mean-reverting, but it still satisfies
\sourceref{eq:SV_no_jump}.  It therefore captures persistent stochastic
variance, but not abrupt short-horizon repricing.

The natural extension is to combine the stochastic variance mechanism
with the jump mechanism developed above.  In a
jump-augmented stochastic-volatility model the price dynamics is replaced
by
\begin{equation}
\frac{dS_t}{S_{t-}}
=
(r-\lambda k)\,dt
+
\sqrt{v_t}\,dW^S_t
+
(e^Y-1)\,dN_t,
\qquad
k:=\mathbb{E}^{\mathbb{Q}}[e^Y-1],
\end{equation}
while the variance continues to follow a stochastic volatility equation,
for instance
\[
dv_t
=
\kappa(\theta-v_t)\,dt+\xi\sqrt{v_t}\,dW^v_t.
\]
This separates the two empirical mechanisms: $v_t$ controls the random
evolution of diffusive variance, whereas the jump component controls
abrupt price discontinuities.  The Bates model is the canonical affine
example of this synthesis, combining the Heston square-root variance
process with Merton-type jumps in the asset price.

The two structural questions left open by the present chapter, namely
the well-posedness of the variance dynamics \sourceref{eq:heston_v} and the
explicit availability of the affine characteristic function
\sourceref{eq:heston_CF_affine}, are the object of the following section.

% ---- next approved source slice ----

\section{From Merton to Heston and Bates}

The preceding sections have systematically disassembled the classical Black-Scholes-Merton framework, isolating its structural deficiencies and introducing the mathematical machinery required to address them. The compound-Poisson and broader L\'evy blocks capture discontinuous repricing. The affine stochastic-volatility block then uses a Cox--Ingersoll--Ross square-root diffusion to model persistent, mean-reverting variance.

This chapter serves as the theoretical synthesis of the report. We comparatively analyze the three canonical continuous-time models that form the foundation of modern quantitative finance: the Merton jump-diffusion model, the Heston pure stochastic volatility model, and finally, the Bates model, which unifies both mechanisms. By examining their distinct characteristic functions and the resulting topologies of their implied volatility surfaces, we demonstrate why a unified framework is mathematically and empirically necessary to capture both short-maturity and long-maturity option pricing anomalies.
\subsection{Merton as Jump-Diffusion Benchmark}

The Merton (1976) model represents the foundational benchmark for jump-diffusion dynamics. It deliberately retains the constant volatility assumption of the classical framework but enriches the continuous geometric Brownian motion by superposing a compound Poisson process. This augmentation mathematically generates heavy-tailed return distributions and introduces realistic, discontinuous price shocks.

Under the risk-neutral pricing measure $\mathbb{Q}$, the asset price process $S_t$ in the Merton framework is governed by the following stochastic differential equation:
\begin{equation}
    \frac{dS_t}{S_{t^-}} = (r - \lambda k) dt + \sigma dW_t + (e^J - 1) dN_t
\end{equation}
where $W_t$ is a standard Brownian motion, $N_t$ is a Poisson process with constant intensity $\lambda > 0$, and $\sigma$ remains a deterministic constant diffusion coefficient. The random variable $J$ represents the discrete jump size in the logarithmic price. Merton specifically assumed that these log-jumps are independent and normally distributed, such that $J \sim \mathcal{N}(\alpha_J, \delta_J^2)$. 

To ensure that the discounted asset price remains a strict martingale under $\mathbb{Q}$, the drift is compensated by the term $\lambda k$, where $k$ is the expected relative price jump:
\begin{equation}
    k = \mathbb{E}^{\mathbb{Q}}[e^J - 1] = e^{\alpha_J + \frac{1}{2}\delta_J^2} - 1
\end{equation}

The structural power of the Merton model is most evident in the frequency domain. By applying the L\'evy-Khintchine representation derived in Chapter 5, the model separates cleanly into a purely continuous Gaussian diffusion and a pure-jump component. Let $X_T = \ln(S_T)$ be the terminal log-price. Its characteristic function takes the exponential form $\varphi_T(u) = \exp(i u X_0 + T \psi_{Merton}(u))$, where the characteristic exponent is exactly:
\begin{equation}
    \psi_{Merton}(u) = iu \left( r - \frac{1}{2}\sigma^2 - \lambda k \right) - \frac{1}{2}\sigma^2 u^2 + \lambda \left( e^{iu\alpha_J - \frac{1}{2}\delta_J^2 u^2} - 1 \right)
\end{equation}

From an empirical pricing perspective, the Merton model substantially improves
the representation of short-maturity tail risk relative to pure diffusion
models. Because the jump component $(e^J-1)dN_t$ allows for discontinuous price
moves, the model can assign non-negligible probability to abrupt repricings even
over short horizons, where the continuous Brownian variance has little time to
accumulate. This mechanism helps generate steep short-maturity smiles and skews,
especially when jump sizes are asymmetric.

However, the Merton framework has an important limitation at longer horizons.
For finite-variance independent jumps, the cumulative jump component becomes
increasingly Gaussian after standardization. More precisely, both the diffusive
variance and the jump variance grow linearly with maturity, while standardized
higher moments such as skewness and excess kurtosis decrease with the horizon.
Consequently, the slope of the implied-volatility skew generated by a pure
jump--diffusion tends to decay at long maturities. This makes the model less
suited to reproducing persistent long-maturity volatility skews without adding a
separate stochastic-volatility mechanism.
\subsection{Heston as Stochastic Volatility Benchmark}

While the Merton model addresses price discontinuities, the Heston (1993) framework serves as the canonical benchmark for capturing the dynamic and persistent nature of market volatility. By relaxing the Black-Scholes assumption of deterministic variance, Heston promotes $v_t$ to a stochastic state variable, effectively modeling "volatility clustering" and the asymmetric feedback between price and variance shocks.

Under the risk-neutral measure $\mathbb{Q}$, the Heston model is defined by the coupled system of SDEs already analyzed in Chapters 7 and 8:
\begin{equation}
    \begin{aligned}
        dS_t &= r S_t dt + \sqrt{v_t} S_t dW_t^S \\
        dv_t &= \kappa(\theta - v_t) dt + \xi \sqrt{v_t} dW_t^v
    \end{aligned}
\end{equation}
with the critical infinitesimal correlation $d\langle W^S, W^v \rangle_t = \rho dt$. 

The primary mechanism for generating the implied volatility skew in this framework is the leverage effect parameter $\rho$. For equity markets, a negative correlation ($\rho < 0$) mathematically implies that a decrease in the asset price is statistically coupled with an increase in instantaneous variance. This structural dependence shifts the probability mass toward the left tail of the return distribution, naturally increasing the price of out-of-the-money (OTM) Put options and generating a downward-sloping volatility skew.

As derived in Equation \sourceref{eq:affine_cf_form}, the affine structure of the CIR variance process allows the characteristic function of the terminal log-price $X_T$ to be expressed as:
\begin{equation}
    \varphi_T(u) = \exp \left( C(\tau, u) + D(\tau, u) v_t + iu X_t \right)
\end{equation}
where the coefficients $C$ and $D$ are the closed-form solutions to the Riccati system.

The empirical advantage of the Heston model lies in its ability to reproduce the persistence of the volatility smile across longer maturities. Because the variance process is mean-reverting and diffusive, the uncertainty regarding future volatility accumulates over time. This ensures that the implied volatility surface remains non-flat even for long-dated options, providing a much better fit than Merton's jump-diffusion for maturities exceeding one year.

However, the Heston model exhibits a structural short-maturity limitation. Because the sample paths of both $S_t$ and $v_t$ are almost surely continuous, the short-time behaviour of the log-price remains locally Gaussian. At the at-the-money level, the implied volatility converges to the spot volatility: \[ \lim_{\tau\to 0} \sigma_{\mathrm{imp}}(K=S_t,\tau) = \sqrt{v_t}. \] Away from the money, the small-maturity asymptotics of implied volatility is more delicate, but the structural message remains the same: a continuous stochastic-volatility diffusion cannot generate order-one discontinuous returns over an arbitrarily short horizon. Therefore, very steep short-dated smiles and skews can be difficult to reproduce without extreme parameters or an explicit jump component. This is precisely the region in which jump--diffusion models provide a complementary mechanism.
\subsection{Bates as a Unified Jump-Volatility Model}

The comparative analysis of the Merton and Heston models reveals a distinct temporal dichotomy: pure jump-diffusions excel at pricing short-term tail risk but fail asymptotically, whereas pure stochastic volatility models effectively capture long-term smile persistence but fail to generate sufficient short-term skew. The natural mathematical evolution, pioneered by David Bates (1996), is to unify these two mechanisms into a single, comprehensive Affine Jump-Diffusion (AJD) framework.

The Bates model explicitly embeds Merton's lognormal jump process into Heston's stochastic volatility environment. Under the risk-neutral measure $\mathbb{Q}$, the joint dynamics of the asset price $S_t$ and its instantaneous variance $v_t$ are governed by the following system of stochastic differential equations:
\begin{equation}
    \begin{aligned}
        \frac{dS_t}{S_{t^-}} &= (r - \lambda k) dt + \sqrt{v_t} dW_t^S + (e^J - 1) dN_t \\
        dv_t &= \kappa(\theta - v_t) dt + \xi \sqrt{v_t} dW_t^v
    \end{aligned}
\end{equation}
with the standard Heston Brownian correlation $d\langle W^S, W^v \rangle_t = \rho dt$. As in the Merton framework, $N_t$ is a Poisson process with intensity $\lambda$, $J \sim \mathcal{N}(\alpha_J, \delta_J^2)$ governs the random jump sizes in the logarithmic price, and $k = \mathbb{E}^{\mathbb{Q}}[e^J - 1]$ is the necessary martingale compensator.

A crucial structural assumption in the standard Bates model is the mutual independence between the jump components ($N_t$, $J$) and the continuous diffusion components ($W_t^S$, $W_t^v$). This orthogonality implies that the arrival of a market crash is purely exogenous and does not instantaneously alter the diffusion variance $v_t$, nor does the level of $v_t$ dictate the frequency of the jumps. 

This independence assumption yields a monumental analytical advantage in the frequency domain. Because the continuous and discontinuous innovations are decoupled, the characteristic exponent of the Bates model is strictly additive. The characteristic function of the terminal log-price $X_T = \ln(S_T)$ simply factors into the product of the Heston characteristic function and the Merton pure-jump characteristic function.

Exploiting the affine structure formalized in Chapter 8, the Bates characteristic function takes the exact closed-form representation:
\begin{equation}
    \varphi_{Bates}(u) = \exp \left( C(\tau, u) + D(\tau, u) v_t + iu X_t + \lambda \tau \left( e^{iu\alpha_J - \frac{1}{2}\delta_J^2 u^2} - 1 - iu k \right) \right)
\end{equation}
where $\tau = T - t$, and the functions $C(\tau, u)$ and $D(\tau, u)$ are the exact identical solutions to the Heston Riccati ODEs derived in Equation \sourceref{eq:riccati_odes}.

Equation \sourceref{eq:bates_cf} encapsulates the ultimate pricing engine of modern quantitative finance. By preserving the exponentially affine form, the Bates model remains fully compatible with the Carr-Madan Fourier inversion formulas \sourceref{eq:carr_madan_pricing}. 

Empirically, the Bates framework is designed to combine the strengths of the two benchmark mechanisms. The jump parameters $(\lambda,\alpha_J,\delta_J)$ primarily affect short-maturity tail risk and the steepness of short-dated smiles, while the stochastic-volatility parameters $(\kappa,\theta,\xi,\rho)$ control the persistence, curvature, and leverage-induced skew of the volatility surface at longer maturities. In this sense, Bates provides a unified affine framework in which jumps and stochastic variance play complementary roles across the maturity dimension.
\subsection{Theoretical Comparison of the Main Mechanisms}

The transition from Merton to Heston, and ultimately to Bates, represents a strategic trade-off between local and asymptotic non-normality. To conclude this synthesis, we evaluate the distinct ways in which these models generate the implied volatility skew and how they resolve the "term structure" problem.
\subsubsection{Probability Tails and the Nature of Non-Normality}

The fundamental difference lies in the mechanism of tail generation. In the Merton jump-diffusion model, non-normality is "injected" into the return distribution through the jump-size distribution $J$. This results in a distribution that is a weighted mixture of Gaussians, where the kurtosis is primarily driven by the jump intensity $\lambda$.

Conversely, in the Heston model, the log-return distribution is a mixture of Gaussians indexed by the random integrated variance $I_T = \int_0^T v_s ds$. Here, the excess kurtosis is generated by the "vol-of-vol" parameter $\xi$, which measures the uncertainty of the variance path. While Merton assumes that shocks are discrete and rare, Heston assumes that volatility itself is a continuous but unpredictable process.

\subsubsection{The Term Structure of the Skew}

The most critical theoretical distinction concerns how the implied volatility skew behaves as the time to maturity $\tau = T - t$ increases. We can formalize the "slope" of the skew near the money as $d \sigma_{imp} / d \ln K$.

\begin{enumerate}
    \item \textbf{Merton's Short-Maturity Dominance:} In the Merton model, the skew is extremely steep for short-dated options. However, as established by the Central Limit Theorem, the jump component's contribution to the total variance scales with $\lambda \tau$, while the diffusion variance scales with $\sigma^2 \tau$. As $\tau \to \infty$, the jump risk is diversified away, and the skew decays at a rate of $O(1/\sqrt{\tau})$.
    \item \textbf{Heston's Long-Maturity Persistence:} In the Heston model, the correlation $\rho$ requires time to impact the terminal distribution. For very short maturities, the variance process is "quasi-deterministic," and the skew is negligible. However, for long maturities, the mean-reversion $\kappa$ and vol-of-vol $\xi$ ensure that the skew remains persistent and does not vanish as rapidly as in the pure jump case.
\end{enumerate}
\subsubsection{Why Bates Provides a Unified Benchmark}

The Bates model overcomes these individual failures by assigning different mathematical "responsibilities" to different parameters. By inspecting the joint dynamics in \sourceref{eq:bates_dynamics}, we can conclude that:
\begin{itemize}[label = --]
    \item The \textbf{Jump parameters} $(\lambda, \alpha_J, \delta_J)$ act as the primary drivers of the short-term skew, capturing the market's "crash-o-phobia" for immediate horizons.
    \item The \textbf{Stochastic Volatility parameters} $(\kappa, \theta, \xi, \rho)$ act as the primary drivers of the long-term term structure, capturing the persistent leverage effect and volatility clustering.
\end{itemize}

As a result, among the three benchmark models considered in this chapter, Bates
is the most flexible framework for representing both short-maturity and
long-maturity features of the implied-volatility surface with a single
consistent parameter set. Its theoretical robustness makes it a standard affine
benchmark for derivative-pricing applications, while more elaborate extensions
such as local-stochastic volatility, rough volatility, or stochastic-intensity
jump models may be required for more demanding empirical calibration tasks.


% Asset/code block omitted pending provenance acceptance.

% Asset/code block omitted pending provenance acceptance.

\section{From Independent Jumps to Hawkes Self-Excitation}

The Bates model still assumes that jump arrivals are independent and governed
by a constant intensity. For an asset exposed to clusters of macroeconomic and
geopolitical repricing, that assumption can be relaxed by making the arrival
intensity itself an adapted state variable. With an exponential Hawkes kernel,
one convenient Markovian representation is
\[
 d\lambda_t=\beta(\bar\lambda-\lambda_t)\,dt+\alpha\,dN_t,
 \qquad \lambda_t>0,
\]
so each event raises near-term intensity by $\alpha$ and the effect decays at
rate $\beta$ toward the exogenous baseline $\bar\lambda$. The branching ratio
$n=\alpha/\beta$ measures endogenous amplification. The stationary regime
requires $0\le n<1$, with stationary mean intensity
$\bar\lambda/(1-n)$ under this parameterisation.

Replacing the constant Poisson intensity in Bates by $\lambda_t$ yields the
Full Bates--Hawkes layer used in this study:
\[
 \frac{dS_t}{S_{t^-}}=(r-\lambda_t k)dt+\sqrt{v_t}\,dW_t^S
 +(e^J-1)dN_t,
 \qquad
 dv_t=\kappa(\theta-v_t)dt+\xi\sqrt{v_t}\,dW_t^v.
\]
The process keeps the Heston variance state and lognormal jump marks while
allowing jump probability to react to past arrivals. With the exponential
kernel, the enlarged state remains affine and can be handled through a
Riccati system and characteristic-function pricing. A rough or power-law
kernel would require a different, generally non-Markovian construction.

Self-excitation is a statistical mechanism for event clustering, not a direct
measurement of news. The filtration used here contains price- and
option-derived information but no classified macroeconomic headlines,
geopolitical events or supply-demand reports. Chapter~9 tests whether the extra
state improves the observed option-surface fit and whether it survives
out-of-sample comparison; future work may augment the filtration with
text-based or other heterogeneous data.

\section{The Theoretical Model Ladder}

Black--Scholes leads to Heston, then Bates with Merton jumps, then Full
Bates--Hawkes self-excitation. Theory closes here; calibration, out-of-sample
tests and simulation follow in Chapter~9.
